\begin{register}{H}{\textcolor{red}{SPI\_CMD\_REG}}{0x{}0000}\label{regdesc:SPICMDREG}
\regfield{(reserved)}{7}{25}{0000000}%
\regfield{SPI\_USR}{1}{24}{{0}}%
\regfield{SPI\_UPDATE}{1}{23}{{0}}%
\regfield{(reserved)}{5}{18}{00000}%
\regfield{(reserved)|\textcolor{red}{SPI\_CONF\_BITLEN}}{18}{0}{{0}}%
\reglabel{Reset}\regnewline%
\begin{regdesc}\begin{reglist}
\label{fielddesc:SPICONFBITLEN}\item [\textcolor{red}{SPI\_CONF\_BITLEN（仅对 SPI2 有效）}] 定义 SPI CONF 阶段的 \hyperref[clk:APBCLK]{APB\_CLK} 周期。可在 CONF 阶段配置。 (R/W)
\label{fielddesc:SPIUPDATE}\item [SPI\_UPDATE] 置位此位，将 SPI 寄存器从 APB 时钟域同步到 SPI 模块时钟域。该位仅用于 SPI 主机模式。 (WT)
\label{fielddesc:SPIUSR}\item [SPI\_USR] 使能用户自定义命令。置位此位将触发一次 SPI 操作。操作完成后，此位自动清零。1：使能；0：禁用。CONF\_buf 不可更改该配置。 (R/W/SC)
\end{reglist}\end{regdesc}
\end{register}


\begin{register}{H}{SPI\_ADDR\_REG}{0x{}0004}\label{regdesc:SPIADDRREG}
\regfield{SPI\_USR\_ADDR\_VALUE}{32}{0}{{0}}%
\reglabel{Reset}\regnewline%
\begin{regdesc}\begin{reglist}
\label{fielddesc:SPIUSRADDRVALUE}\item [SPI\_USR\_ADDR\_VALUE] 从机地址。可在 CONF 阶段配置。 (R/W)
\end{reglist}\end{regdesc}
\end{register}


\begin{register}{H}{\textcolor{red}{SPI\_USER\_REG}}{0x{}0010}\label{regdesc:SPIUSERREG}
\regfield{SPI\_USR\_COMMAND}{1}{31}{{1}}%
\regfield{SPI\_USR\_ADDR}{1}{30}{{0}}%
\regfield{SPI\_USR\_DUMMY}{1}{29}{{0}}%
\regfield{SPI\_USR\_MISO}{1}{28}{{0}}%
\regfield{SPI\_USR\_MOSI}{1}{27}{{0}}%
\regfield{SPI\_USR\_DUMMY\_IDLE}{1}{26}{{0}}%
\regfield{SPI\_USR\_MOSI\_HIGHPART}{1}{25}{{0}}%
\regfield{SPI\_USR\_MISO\_HIGHPART}{1}{24}{{0}}%
\regfield{(reserved)}{6}{18}{000000}%
\regfield{SPI\_SIO}{1}{17}{{0}}%
\regfield{(reserved)}{1}{16}{0}%
\regfield{(reserved)|\textcolor{red}{SPI\_USR\_CONF\_NXT}}{1}{15}{{0}}%
\regfield{(reserved)|\textcolor{red}{SPI\_FWRITE\_OCT}}{1}{14}{{0}}%
\regfield{SPI\_FWRITE\_QUAD}{1}{13}{{0}}%
\regfield{SPI\_FWRITE\_DUAL}{1}{12}{{0}}%
\regfield{(reserved)}{2}{10}{00}%
\regfield{SPI\_CK\_OUT\_EDGE}{1}{9}{{0}}%
\regfield{SPI\_RSCK\_I\_EDGE}{1}{8}{{0}}%
\regfield{SPI\_CS\_SETUP}{1}{7}{{1}}%
\regfield{SPI\_CS\_HOLD}{1}{6}{{1}}%
\regfield{SPI\_TSCK\_I\_EDGE}{1}{5}{{0}}%
\regfield{(reserved)|\textcolor{red}{SPI\_OPI\_MODE}}{1}{4}{{0}}%
\regfield{SPI\_QPI\_MODE}{1}{3}{{0}}%
\regfield{(reserved)}{2}{1}{00}%
\regfield{SPI\_DOUTDIN}{1}{0}{{0}}%
\reglabel{Reset}\regnewline%
\begin{regdesc}\begin{reglist}
\label{fielddesc:SPIDOUTDIN}\item [SPI\_DOUTDIN] 使能全双工模式。1：使能；0：禁用。可在 CONF 阶段配置。 (R/W)
\label{fielddesc:SPIQPIMODE}\item [SPI\_QPI\_MODE] 1：使能 QPI 模式。0：禁用 QPI 模式。SPI 主机模式和从机模式均支持该配置。可在 CONF 阶段配置。 (R/W/SS/SC)
\label{fielddesc:SPIOPIMODE}\item [\textcolor{red}{SPI\_OPI\_MODE（仅对 SPI2 有效）}] 1：使能 OPI 模式，则所有阶段均采用 8-bit 模式。0：禁用 OPI 模式。仅有 SPI 主机模式支持该配置。可在 CONF 阶段配置。 (R/W)
\label{fielddesc:SPITSCKIEDGE}\item [SPI\_TSCK\_I\_EDGE] 在从机模式下，此位可用于配置 SPI 四种时钟模式，见章节 \ref{GP-SPI slave clock control}。 (R/W)
\label{fielddesc:SPICSHOLD}\item [SPI\_CS\_HOLD] SPI 处于完成 (DONE) 阶段时，SPI CS 保持低电平。1：使能此功能；0：禁用此功能。可在 CONF 阶段配置。 (R/W)
\label{fielddesc:SPICSSETUP}\item [SPI\_CS\_SETUP] SPI 处于准备 (PREP) 阶段时，使能 SPI CS。1：启用此功能；0：禁用此功能。可在 CONF 阶段配置。 (R/W)
\label{fielddesc:SPIRSCKIEDGE}\item [SPI\_RSCK\_I\_EDGE] 在从机模式下，此位可用于配置四种 SPI 时钟模式，见章节 \ref{GP-SPI slave clock control}。 (R/W)
\label{fielddesc:SPICKOUTEDGE}\item [SPI\_CK\_OUT\_EDGE] 该位与 SPI\_MOSI\_DELAY\_MODE 一起用于控制 MOSI 信号延迟模式。可在 CONF 阶段配置。 (R/W)
\label{fielddesc:SPIFWRITEDUAL}\item [SPI\_FWRITE\_DUAL] 在写操作 (DOUT) 阶段，读数据的方式为 2-bit 方式。可在 CONF 阶段配置。 (R/W)
\label{fielddesc:SPIFWRITEQUAD}\item [SPI\_FWRITE\_QUAD] 在写操作 (DOUT) 阶段，读数据的方式为 4-bit 方式。可在 CONF 阶段配置。 (R/W)
\label{fielddesc:SPIFWRITEOCT}\item [\textcolor{red}{SPI\_FWRITE\_OCT（仅对 SPI2 有效）}] 在写操作 (DOUT) 阶段，读数据的方式为 8-bit 方式。可在 CONF 阶段配置。 (R/W)
\label{fielddesc:SPIUSRCONFNXT}\item [\textcolor{red}{SPI\_USR\_CONF\_NXT（仅对 SPI2 有效）}]
使能下一次传输事务的 CONF 阶段。可在 CONF 阶段配置。 (R/W)
\begin{itemize}
    \item 置位此位，则本次分段配置传输继续进行，开始下一次传输事务。
    \item 清除此位，则当前传输事务结束后，本次分段配置传输结束。或者，当前的传输模式不是分段配置传输。
\end{itemize}
\label{fielddesc:SPISIO}\item [SPI\_SIO] 配置三线半双工通信。MOSI 和 MISO 信号共用一个管脚。1：使能三线半双工通信；0：禁用三线半双工通信。可在 CONF 阶段配置。 (R/W)
\label{fielddesc:SPIUSRMISOHIGHPART}\item [SPI\_USR\_MISO\_HIGHPART] 在读数据阶段，仅访问高位 buffer：\hyperref[regdesc:SPIW8REG]{SPI\_W8\_REG} \verb+~+ \hyperref[regdesc:SPIW15REG]{SPI\_W15\_REG}。1：使能；0：禁用。可在 CONF 阶段配置。 (R/W)

\item [见下页]
\end{reglist}\end{regdesc}
\end{register}

\addtocounter{Regfloat}{-1}
\begin{register}{H}{\textcolor{red}{SPI\_USER\_REG}}{0x{}0010}
\begin{regdesc}\begin{reglist}
\item [接上页]

\label{fielddesc:SPIUSRMOSIHIGHPART}\item [SPI\_USR\_MOSI\_HIGHPART] 在写数据阶段，仅访问高位 buffer：\hyperref[regdesc:SPIW8REG]{SPI\_W8\_REG} \verb+~+ \hyperref[regdesc:SPIW15REG]{SPI\_W15\_REG}。1：使能；0：禁用。可在 CONF 阶段配置。 (R/W)
\label{fielddesc:SPIUSRDUMMYIDLE}\item [SPI\_USR\_DUMMY\_IDLE] 置位此位，在 DUMMY 阶段禁用 SPI 时钟。可在 CONF 阶段配置。 (R/W)
\label{fielddesc:SPIUSRMOSI}\item [SPI\_USR\_MOSI] 置位此位，使能一次操作的写数据 (DOUT) 阶段。可在 CONF 阶段配置。 (R/W)
\label{fielddesc:SPIUSRMISO}\item [SPI\_USR\_MISO] 置位此位，使能一次操作的读数据 (DIN) 阶段。可在 CONF 阶段配置。 (R/W)
\label{fielddesc:SPIUSRDUMMY}\item [SPI\_USR\_DUMMY] 置位此位，使能一次操作的 DUMMY 阶段。可在 CONF 阶段配置。 (R/W)
\label{fielddesc:SPIUSRADDR}\item [SPI\_USR\_ADDR] 置位此位，使能一次操作的地址 (ADDR) 阶段。可在 CONF 阶段配置。 (R/W)
\label{fielddesc:SPIUSRCOMMAND}\item [SPI\_USR\_COMMAND] 置位此位，使能一次操作的命令 (CMD) 阶段。可在 CONF 阶段配置。 (R/W)
\end{reglist}\end{regdesc}
\end{register}


\begin{register}{H}{SPI\_USER1\_REG}{0x{}0014}\label{regdesc:SPIUSER1REG}
\regfield{SPI\_USR\_ADDR\_BITLEN}{5}{27}{{23}}%
\regfield{SPI\_CS\_HOLD\_TIME}{5}{22}{{0x{}1}}%
\regfield{SPI\_CS\_SETUP\_TIME}{5}{17}{{0}}%
\regfield{SPI\_MST\_WFULL\_ERR\_END\_EN}{1}{16}{{1}}%
\regfield{(reserved)}{8}{8}{00000000}%
\regfield{SPI\_USR\_DUMMY\_CYCLELEN}{8}{0}{{7}}%
\reglabel{Reset}\regnewline%
\begin{regdesc}\begin{reglist}
\label{fielddesc:SPIUSRDUMMYCYCLELEN}\item [SPI\_USR\_DUMMY\_CYCLELEN] DUMMY 阶段的时长，单位：SPI\_CLK 时钟周期。寄存器值为（实际需要的周期数 - 1）。可在 CONF 阶段配置。 (R/W)
\label{fielddesc:SPIMSTWFULLERRENDEN}\item [SPI\_MST\_WFULL\_ERR\_END\_EN] 1：在 GP-SPI 主机全双工或半双工模式下，如果发生 SPI RX AFIFO 满错误，则 SPI 传输将终止。0：在 GP-SPI 主机全双工或半双工模式下，如果发生 SPI RX AFIFO 满错误，SPI 传输将不被终止。 (R/W)
\label{fielddesc:SPICSSETUPTIME}\item [SPI\_CS\_SETUP\_TIME] 准备 (PREP) 阶段的时长，单位：SPI\_CLK 时钟周期。此值等于预期周期数 - 1。此字段与 \hyperref[fielddesc:SPICSSETUP]{SPI\_CS\_SETUP} 搭配使用。可在 CONF 阶段配置。 (R/W)
\label{fielddesc:SPICSHOLDTIME}\item [SPI\_CS\_HOLD\_TIME] CS 的延迟周期。单位：SPI\_CLK 时钟周期。此字段与 \hyperref[fielddesc:SPICSHOLD]{SPI\_CS\_HOLD} 搭配使用。可在 CONF 阶段配置。 (R/W)
\label{fielddesc:SPIUSRADDRBITLEN}\item [SPI\_USR\_ADDR\_BITLEN] 地址阶段的位长。此值为（预期位数 - 1）。可在 CONF 阶段配置。 (R/W)
\end{reglist}\end{regdesc}
\end{register}


\begin{register}{H}{SPI\_USER2\_REG}{0x{}0018}\label{regdesc:SPIUSER2REG}
\regfield{SPI\_USR\_COMMAND\_BITLEN}{4}{28}{{7}}%
\regfield{SPI\_MST\_REMPTY\_ERR\_END\_EN}{1}{27}{{1}}%
\regfield{(reserved)}{11}{16}{00000000000}%
\regfield{SPI\_USR\_COMMAND\_VALUE}{16}{0}{{0}}%
\reglabel{Reset}\regnewline%
\begin{regdesc}\begin{reglist}
\label{fielddesc:SPIUSRCOMMANDVALUE}\item [SPI\_USR\_COMMAND\_VALUE] 命令值。可在 CONF 阶段配置。 (R/W)
\label{fielddesc:SPIMSTREMPTYERRENDEN}\item [SPI\_MST\_REMPTY\_ERR\_END\_EN] 1：在 GP-SPI 主机全双工或半双工模式下，如果发生 SPI TX AFIFO 空错误，则 SPI 传输将终止。0：在 GP-SPI 主机全双工或半双工模式下，如果发生 SPI TX AFIFO 空错误，则 SPI 传输将不会被终止。 (R/W)
\label{fielddesc:SPIUSRCOMMANDBITLEN}\item [SPI\_USR\_COMMAND\_BITLEN] 命令阶段的位长。此值为（预期位长 - 1）。可在 CONF 阶段配置。 (R/W)
\end{reglist}\end{regdesc}
\end{register}


\begin{register}{H}{\textcolor{red}{SPI\_CTRL\_REG}}{0x{}0008}\label{regdesc:SPICTRLREG}
\regfield{(reserved)}{5}{27}{00000}%
\regfield{SPI\_WR\_BIT\_ORDER}{2}{25}{{0}}%
\regfield{SPI\_RD\_BIT\_ORDER}{2}{23}{{0}}%
\regfield{(reserved)}{1}{22}{0}%
\regfield{SPI\_WP\_POL}{1}{21}{{1}}%
\regfield{SPI\_HOLD\_POL}{1}{20}{{1}}%
\regfield{SPI\_D\_POL}{1}{19}{{1}}%
\regfield{SPI\_Q\_POL}{1}{18}{{1}}%
\regfield{(reserved)}{1}{17}{0}%
\regfield{(reserved)|\textcolor{red}{SPI\_FREAD\_OCT}}{1}{16}{{0}}%
\regfield{SPI\_FREAD\_QUAD}{1}{15}{{0}}%
\regfield{SPI\_FREAD\_DUAL}{1}{14}{{0}}%
\regfield{(reserved)}{3}{11}{000}%
\regfield{(reserved)|\textcolor{red}{SPI\_FCMD\_OCT}}{1}{10}{{0}}%
\regfield{SPI\_FCMD\_QUAD}{1}{9}{{0}}%
\regfield{SPI\_FCMD\_DUAL}{1}{8}{{0}}%
\regfield{(reserved)|\textcolor{red}{SPI\_FADDR\_OCT}}{1}{7}{{0}}%
\regfield{SPI\_FADDR\_QUAD}{1}{6}{{0}}%
\regfield{SPI\_FADDR\_DUAL}{1}{5}{{0}}%
\regfield{(reserved)}{1}{4}{0}%
\regfield{SPI\_DUMMY\_OUT}{1}{3}{{0}}%
\regfield{(reserved)}{3}{0}{000}%
\reglabel{Reset}\regnewline%
\begin{regdesc}\begin{reglist}
\label{fielddesc:SPIDUMMYOUT}\item [SPI\_DUMMY\_OUT] 可在 CONF 阶段配置。 (R/W)
\begin{itemize}
    \item \textbf{SPI2}
    \begin{itemize}
        \item 0：在 DUMMY 阶段，不输出 FSPI 总线信号。
        \item 1：在 DUMMY 阶段，输出 FSPI 总线信号。
    \end{itemize}
    \item \textbf{SPI3}：表示在 DUMMY 阶段，SPI3 输出的信号电平。
\end{itemize}

\label{fielddesc:SPIFADDRDUAL}\item [SPI\_FADDR\_DUAL] 在地址 (ADDR) 阶段采用 2-bit 模式。1：使能；0：禁用。可在 CONF 阶段配置。 (R/W)
\label{fielddesc:SPIFADDRQUAD}\item [SPI\_FADDR\_QUAD] 在地址 (ADDR) 阶段采用 4-bit 模式。1：使能；0：禁用。可在 CONF 阶段配置。 (R/W)
\label{fielddesc:SPIFADDROCT}\item [\textcolor{red}{SPI\_FADDR\_OCT（仅对 SPI2 有效）}] 在地址 (ADDR) 阶段采用 8-bit 模式。1：使能；0：禁用。可在 CONF 阶段配置。 (R/W)
\label{fielddesc:SPIFCMDDUAL}\item [SPI\_FCMD\_DUAL] 在命令 (CMD) 阶段采用 2-bit 模式。1：使能；0：禁用。可在 CONF 阶段配置。 (R/W)
\label{fielddesc:SPIFCMDQUAD}\item [SPI\_FCMD\_QUAD] 在命令 (CMD) 阶段采用 4-bit 模式。1：使能；0：禁用。可在 CONF 阶段配置。 (R/W)
\label{fielddesc:SPIFCMDOCT}\item [\textcolor{red}{SPI\_FCMD\_OCT (for SP2 only)}] 在命令 (CMD) 阶段采用 8-bit 模式。1：使能；0：禁用。可在 CONF 阶段配置。 (R/W)
\label{fielddesc:SPIFREADDUAL}\item [SPI\_FREAD\_DUAL] 在读数据 (DIN) 阶段采用 2-bit 模式。1：使能；0：禁用。可在 CONF 阶段配置。 (R/W)
\label{fielddesc:SPIFREADQUAD}\item [SPI\_FREAD\_QUAD] 在读数据 (DIN) 阶段采用 4-bit 模式。1：使能；0：禁用。可在 CONF 阶段配置。 (R/W)
\label{fielddesc:SPIFREADOCT}\item [\textcolor{red}{SPI\_FREAD\_OCT (for SP2 only)}] 在读数据 (DIN) 阶段采用 8-bit 模式。1：使能；0：禁用。可在 CONF 阶段配置。 (R/W)
\label{fielddesc:SPIQPOL}\item [SPI\_Q\_POL] 此位用于设置 MISO 的极性。1：高；0：低。可在 CONF 阶段配置。 (R/W)
\label{fielddesc:SPIDPOL}\item [SPI\_D\_POL] 此位用于设置 MOSI 的极性。1：高；0：低。可在 CONF 阶段配置。 (R/W)

\item [见下页]
\end{reglist}\end{regdesc}
\end{register}

\addtocounter{Regfloat}{-1}
\begin{register}{H}{\textcolor{red}{SPI\_CTRL\_REG}}{0x{}0008}
\begin{regdesc}\begin{reglist}
\item [接上页]

\label{fielddesc:SPIHOLDPOL}\item [SPI\_HOLD\_POL] 此位用于设置在 SPI 空闲状态下，SPI\_HOLD 的输出值。1：输出高电平；0：输出低电平。可在 CONF 阶段配置。 (R/W)
\label{fielddesc:SPIWPPOL}\item [SPI\_WP\_POL] 此位用于设置在 SPI 空闲状态下，WP 信号的输出值。1：输出高电平；0：输出低电平。可在 CONF 阶段配置。 (R/W)
\label{fielddesc:SPIRDBITORDER}\item [SPI\_RD\_BIT\_ORDER] 在读数据 (MISO) 阶段，1：先读低有效位；0：先读高有效位。可在 CONF 阶段配置。 (R/W)
\label{fielddesc:SPIWRBITORDER}\item [SPI\_WR\_BIT\_ORDER] 在命令 (CMD)、地址 (ADDR) 和写数据 (MOSI) 阶段，1：先发送低有效位；0：先发送高有效位。可在 CONF 阶段配置。 (R/W)
\end{reglist}\end{regdesc}
\end{register}


\begin{register}{H}{SPI\_MS\_DLEN\_REG}{0x{}001C}\label{regdesc:SPIMSDLENREG}
\regfield{(reserved)}{14}{18}{00000000000000}%
\regfield{SPI\_MS\_DATA\_BITLEN}{18}{0}{{0}}%
\reglabel{Reset}\regnewline%
\begin{regdesc}\begin{reglist}
\label{fielddesc:SPIMSDATABITLEN}\item [SPI\_MS\_DATA\_BITLEN] 该字段用于配置主机模式下 DMA 控制或 CPU 控制的 SPI 传输的数据位长。也可用于配置从机模式下 DMA 控制的传输中接收数据的位长。该值等于需要的位长 - 1。可在 CONF 阶段配置。 (R/W)
\end{reglist}\end{regdesc}
\end{register}


\begin{register}{H}{\textcolor{red}{SPI\_MISC\_REG（仅对 SPI2 有效）}}{0x{}0020}\label{regdesc:SPIMISCREG}
\regfield{SPI\_QUAD\_DIN\_PIN\_SWAP}{1}{31}{{0}}%
\regfield{SPI\_CS\_KEEP\_ACTIVE}{1}{30}{{0}}%
\regfield{SPI\_CK\_IDLE\_EDGE}{1}{29}{{0}}%
\regfield{(reserved)}{4}{25}{0000}%
\regfield{SPI\_DQS\_IDLE\_EDGE}{1}{24}{{0}}%
\regfield{SPI\_SLAVE\_CS\_POL}{1}{23}{{0}}%
\regfield{(reserved)}{3}{20}{000}%
\regfield{SPI\_CMD\_DTR\_EN}{1}{19}{{0}}%
\regfield{SPI\_ADDR\_DTR\_EN}{1}{18}{{0}}%
\regfield{SPI\_DATA\_DTR\_EN}{1}{17}{{0}}%
\regfield{SPI\_CLK\_DATA\_DTR\_EN}{1}{16}{{0}}%
\regfield{(reserved)}{3}{13}{000}%
\regfield{SPI\_MASTER\_CS\_POL}{6}{7}{{0}}%
\regfield{SPI\_CK\_DIS}{1}{6}{{0}}%
\regfield{SPI\_CS5\_DIS}{1}{5}{{1}}%
\regfield{SPI\_CS4\_DIS}{1}{4}{{1}}%
\regfield{SPI\_CS3\_DIS}{1}{3}{{1}}%
\regfield{SPI\_CS2\_DIS}{1}{2}{{1}}%
\regfield{SPI\_CS1\_DIS}{1}{1}{{1}}%
\regfield{SPI\_CS0\_DIS}{1}{0}{{0}}%
\reglabel{Reset}\regnewline%
\begin{regdesc}\begin{reglist}
\label{fielddesc:SPICS0DIS}\item [SPI\_CS0\_DIS] SPI CS0 管脚使能。1：禁用 CS0；0：SPI CS0 信号来自 CS0 管脚或输出至 CS0 管脚。可在 CONF 阶段配置。 (R/W)
\label{fielddesc:SPICS1DIS}\item [SPI\_CS1\_DIS] SPI CS1 管脚使能。1：禁用 CS1；0：SPI CS1 信号来自 CS1 管脚或输出至 CS1 管脚。可在 CONF 阶段配置。 (R/W)
\label{fielddesc:SPICS2DIS}\item [SPI\_CS2\_DIS] SPI CS2 管脚使能。1：禁用 CS2；0：SPI CS2 信号来自 CS2 管脚或输出至 CS2 管脚。可在 CONF 阶段配置。 (R/W)
\label{fielddesc:SPICS3DIS}\item [SPI\_CS3\_DIS] SPI CS3 管脚使能。1：禁用 CS3；0：SPI CS3 信号来自 CS3 管脚或输出至 CS3 管脚。可在 CONF 阶段配置。 (R/W)
\label{fielddesc:SPICS4DIS}\item [SPI\_CS4\_DIS] SPI CS4 管脚使能。1：禁用 CS4；0：SPI CS4 信号来自 CS4 管脚或输出至 CS4 管脚。可在 CONF 阶段配置。 (R/W)
\label{fielddesc:SPICS5DIS}\item [SPI\_CS5\_DIS] SPI CS5 管脚使能。1：禁用 CS5；0：SPI CS5 信号来自 CS5 管脚或输出至 CS5 管脚。可在 CONF 阶段配置。 (R/W)
\label{fielddesc:SPICKDIS}\item [SPI\_CK\_DIS] 1：停止 SPI\_CLK 输出信号；0：使能 SPI\_CLK 输出信号。可在 CONF 阶段配置。 (R/W)
\label{fielddesc:SPIMASTERCSPOL}\item [SPI\_MASTER\_CS\_POL] 主机模式下，SPI\_MASTER\_CS\_POL[\regindex{i}] 用于配置 SPI CS\regindex{i} 的极性，\regindex{i} = 0 \verb+~+ 5。0：CS\regindex{i} 低电平有效。1：CS\regindex{i} 高电平有效。可在 CONF 阶段配置。 (R/W) 
\label{fielddesc:SPICLKDATADTREN}\item [SPI\_CLK\_DATA\_DTR\_EN] 1：在 SPI 主机模式下，SPI\_CLK、DATA 和 SPI\_DQS 采用 DDR 模式。0：在 SPI 主机模式下，仅 SPI\_DQS 采用 DDR 模式。该位与 bit 17/18/19 一起使用。 (R/W)
\label{fielddesc:SPIDATADTREN}\item [SPI\_DATA\_DTR\_EN] 1：SPI 在主机 1/2/4/8-bit 模式下，DIN 和 DOUT 阶段的时钟和数据均采用 DDR 模式。0：SPI DIN 和 DOUT 阶段的时钟和数据均采用 SDR 模式。可在 CONF 阶段配置。 (R/W)
\label{fielddesc:SPIADDRDTREN}\item [SPI\_ADDR\_DTR\_EN] 1：SPI 在主机 1/2/4/8-bit 模式下，SPI\_SEND\_ADDR 阶段的时钟和数据均采用 DDR 模式。0：SPI\_SEND\_ADDR 阶段的时钟和数据模式均采用 SDR 模式。可在 CONF 阶段配置。 (R/W)
\label{fielddesc:SPICMDDTREN}\item [SPI\_CMD\_DTR\_EN] 1：SPI 在主机 1/2/4/8-bit 模式下，SPI\_SEND\_CMD 阶段的时钟和数据均采用 DDR 模式。0：SPI\_SEND\_CMD 阶段的时钟和数据模式均采用 SDR 模式。可在 CONF 阶段配置。 (R/W)
\label{fielddesc:SPISLAVECSPOL}\item [SPI\_SLAVE\_CS\_POL] 选择 SPI 从机输入信号 CS 的极性。1：反相；0：保持不变。可在 CONF 阶段配置。 (R/W)
\label{fielddesc:SPIDQSIDLEEDGE}\item [SPI\_DQS\_IDLE\_EDGE] SPI\_DQS 的默认值。0：低电平；1：高电平。可在 CONF 阶段配置。 (R/W)
\item [见下页]
\end{reglist}\end{regdesc}
\end{register}

\addtocounter{Regfloat}{-1}
\begin{register}{H}{\textcolor{red}{SPI\_MISC\_REG（仅对 SPI2 有效）}}{0x{}0020}
\begin{regdesc}\begin{reglist}
\item [接上页]
\label{fielddesc:SPICKIDLEEDGE}\item [SPI\_CK\_IDLE\_EDGE] 1：SPI CLK 线在空闲状态时保持高电平；0：SPI CLK 线在空闲状态时保持低电平。可在 CONF 阶段配置。 (R/W)
\label{fielddesc:SPICSKEEPACTIVE}\item [SPI\_CS\_KEEP\_ACTIVE] 置位此位，则 SPI CS 线保持低电平。可在 CONF 阶段配置。 (R/W)
\label{fielddesc:SPIQUADDINPINSWAP}\item [SPI\_QUAD\_DIN\_PIN\_SWAP] 使能 SPI Quad 输入交换。1：FSPID 与 FSPIQ 交换，FSPIWP 与 FSPIHD 交换。0：禁止 SPI Quad 输入交换。可在 CONF 阶段配置。 (R/W)
\end{reglist}\end{regdesc}
\end{register}

\begin{register}{H}{\textcolor{red}{SPI\_MISC\_REG（仅对 SPI3 有效）}}{0x{}0020}\label{regdesc:SPI3MISCREG}
\regfield{SPI\_QUAD\_DIN\_PIN\_SWAP}{1}{31}{{0}}%
\regfield{SPI\_CS\_KEEP\_ACTIVE}{1}{30}{{0}}%
\regfield{SPI\_CK\_IDLE\_EDGE}{1}{29}{{0}}%
\regfield{(reserved)}{5}{24}{00000}%
\regfield{SPI\_SLAVE\_CS\_POL}{1}{23}{{0}}%
\regfield{(reserved)}{13}{10}{0000000000000}%
\regfield{SPI\_MASTER\_CS\_POL}{3}{7}{{0}}%
\regfield{SPI\_CK\_DIS}{1}{6}{{0}}%
\regfield{(reserved)}{3}{3}{000}%
\regfield{SPI\_CS2\_DIS}{1}{2}{{1}}%
\regfield{SPI\_CS1\_DIS}{1}{1}{{1}}%
\regfield{SPI\_CS0\_DIS}{1}{0}{{0}}%
\reglabel{Reset}\regnewline%
\begin{regdesc}\begin{reglist}
\label{fielddesc:SPI3CS0DIS}\item [SPI\_CS0\_DIS] SPI CS0 管脚使能。1：禁用 CS0；0：SPI CS0 信号来自 CS0 管脚或输出至 CS0 管脚。可在 CONF 阶段配置。 (R/W)
\label{fielddesc:SPI3CS1DIS}\item [SPI\_CS1\_DIS] SPI CS1 管脚使能。1：禁用 CS1；0：SPI CS1 信号来自 CS1 管脚或输出至 CS1 管脚。可在 CONF 阶段配置。 (R/W)
\label{fielddesc:SPI3CS2DIS}\item [SPI\_CS2\_DIS] SPI CS2 管脚使能。1：禁用 CS2；0：1：禁用 CS2；0：SPI CS2 信号来自 CS2 管脚或输出至 CS2 管脚。可在 CONF 阶段配置。 (R/W)
\label{fielddesc:SPI3CKDIS}\item [SPI\_CK\_DIS] 1：停止 SPI\_CLK 输出信号；0：使能 SPI\_CLK 输出信号。可在 CONF 阶段配置。 (R/W)
\label{fielddesc:SPI3MASTERCSPOL}\item [SPI\_MASTER\_CS\_POL] 主机模式下，SPI\_MASTER\_CS\_POL[\regindex{i}] 用于配置 SPI CS\regindex{i} 的极性，\regindex{i} = 0 \verb+~+ 2。0：CS\regindex{i} 低电平有效。1：CS\regindex{i} 高电平有效。可在 CONF 阶段配置。 (R/W) 
\label{fielddesc:SPI3SLAVECSPOL}\item [SPI\_SLAVE\_CS\_POL] 选择 SPI 从机输入信号 CS 的极性。1：反相；0：保持不变。可在 CONF 阶段配置。 (R/W)
\label{fielddesc:SPI3CKIDLEEDGE}\item [SPI\_CK\_IDLE\_EDGE] 1：SPI CLK 线在空闲状态时保持高电平；0：SPI CLK 线在空闲状态时保持低电平。可在 CONF 阶段配置。 (R/W)
\label{fielddesc:SPI3CSKEEPACTIVE}\item [SPI\_CS\_KEEP\_ACTIVE] 置位此位，则 SPI CS 线保持低电平。可在 CONF 阶段配置。 (R/W)
\label{fielddesc:SPI3QUADDINPINSWAP}\item [SPI\_QUAD\_DIN\_PIN\_SWAP] 1：使能 SPI Quad 输入交换；0：禁止 SPI Quad 输入交换。可在 CONF 阶段配置。 (R/W)
\end{reglist}\end{regdesc}
\end{register}


\begin{register}{H}{SPI\_DMA\_CONF\_REG}{0x{}0030}\label{regdesc:SPIDMACONFREG}
\regfield{SPI\_DMA\_AFIFO\_RST}{1}{31}{{0}}%
\regfield{SPI\_BUF\_AFIFO\_RST}{1}{30}{{0}}%
\regfield{SPI\_RX\_AFIFO\_RST}{1}{29}{{0}}%
\regfield{SPI\_DMA\_TX\_ENA}{1}{28}{{0}}%
\regfield{SPI\_DMA\_RX\_ENA}{1}{27}{{0}}%
\regfield{(reserved)}{5}{22}{00000}%
\regfield{SPI\_RX\_EOF\_EN}{1}{21}{{0}}%
\regfield{SPI\_SLV\_TX\_SEG\_TRANS\_CLR\_EN}{1}{20}{{0}}%
\regfield{SPI\_SLV\_RX\_SEG\_TRANS\_CLR\_EN}{1}{19}{{0}}%
\regfield{SPI\_DMA\_SLV\_SEG\_TRANS\_EN}{1}{18}{{0}}%
\regfield{(reserved)}{16}{2}{0000000000000000}%
\regfield{SPI\_DMA\_INFIFO\_FULL}{1}{1}{{1}}%
\regfield{SPI\_DMA\_OUTFIFO\_EMPTY}{1}{0}{{1}}%
\reglabel{Reset}\regnewline%
\begin{regdesc}\begin{reglist}
\label{fielddesc:SPIDMAOUTFIFOEMPTY}\item [SPI\_DMA\_OUTFIFO\_EMPTY] 记录 DMA TX FIFO 的状态。1：DMA TX FIFO 尚未就绪，不能发送数据。0：DMA TX FIFO 已就绪，可以发送数据。 (RO)
\label{fielddesc:SPIDMAINFIFOFULL}\item [SPI\_DMA\_INFIFO\_FULL] 记录 DMA RX FIFO 的状态。1：DMA RX FIFO 尚未就绪，不能接收数据。0：DMA RX FIFO 已就绪，可以接收数据。 (RO)
\label{fielddesc:SPIDMASLVSEGTRANSEN}\item [SPI\_DMA\_SLV\_SEG\_TRANS\_EN] 1：使能从机半双工通信方式下，DMA 控制的连续传输模式。0：禁用。 (R/W)
\label{fielddesc:SPISLVRXSEGTRANSCLREN}\item [SPI\_SLV\_RX\_SEG\_TRANS\_CLR\_EN] 在 DMA 控制的半双工从机模式下，当 DMA RX buffer 小于实际接
收的数据长度时，1：后续传输的数据都不接收；0：本次传输的数据不接收，下次传输时, 如果
DMA RX buffer 长度不为零，则继续接收，否则不接收。 (R/W)
\label{fielddesc:SPISLVTXSEGTRANSCLREN}\item [SPI\_SLV\_TX\_SEG\_TRANS\_CLR\_EN] 在 DMA 控制的半双工从机模式下，当 DMA TX buffer 大小小于实际发送的数据长度时，1：后续传输的数据都不更新，发送同一个旧数据；0：本次传输的数据都不更新，下次传输时, 如果 DMA TX FIFO 填充了新的数据，则继续发送新数据，否则发送数据不更新。 (R/W)
\label{fielddesc:SPIRXEOFEN}\item [SPI\_RX\_EOF\_EN] 1：在 DMA 控制的数据传输过程中，如果 DMA 传输的数据比特数等于 (\hyperref[fielddesc:SPIMSDATABITLEN]{SPI\_MS\_DATA\_BITLEN} + 1)，则硬件会置位 GDMA\_IN\_SUC\_EOF\_CH\textcolor{red}{\it{n}}\_INT\_RAW。0：在非分段配置传输模式下，GDMA\_IN\_SUC\_EOF\_CH\textcolor{red}{\it{n}}\_INT\_RAW 由 \hyperref[int:SPITRANSDONEINT]{SPI\_TRANS\_DONE\_INT} 置位；或在分段配置传输模式下，由 \hyperref[int:SPIDMASEGTRANSDONEINT]{SPI\_DMA\_SEG\_TRANS\_DONE\_INT} 置位。(R/W)

\label{fielddesc:SPIDMARXENA}\item [SPI\_DMA\_RX\_ENA] 置位此位，使能 DMA 控制的接收数据模式。 (R/W)
\label{fielddesc:SPIDMATXENA}\item [SPI\_DMA\_TX\_ENA] 置位此位，使能 DMA 控制的发送数据模式。 (R/W)
\label{fielddesc:SPIRXAFIFORST}\item [SPI\_RX\_AFIFO\_RST] 置位此位，复位图 \ref{fig:spi-master-data-flow-control} 和图  \ref{fig:spi-slave-data-flow-control} 中的 spi\_rx\_afifo。spi\_rx\_afifo 将在 SPI 主机和从机传输中用于接收数据。 (WT)

\label{fielddesc:SPIBUFAFIFORST}\item [SPI\_BUF\_AFIFO\_RST] 置位此位，复位图 \ref{fig:spi-master-data-flow-control}  和图  \ref{fig:spi-slave-data-flow-control} 中的 buf\_tx\_afifo。buf\_tx\_afifo 将在 CPU 控制的从机传输或主机传输中用于发送数据。 (WT)

\label{fielddesc:SPIDMAAFIFORST}\item [SPI\_DMA\_AFIFO\_RST] 置位此位，复位图  \ref{fig:spi-slave-data-flow-control} 中的 dma\_tx\_afifo，dma\_tx\_afifo 在 DMA 控制的从机传输中用于发送数据。 (WT)
\end{reglist}\end{regdesc}
\end{register}


\begin{register}{H}{\textcolor{red}{SPI\_SLAVE\_REG}}{0x{}00E0}\label{regdesc:SPISLAVEREG}
\regfield{(reserved)}{3}{29}{000}%
\regfield{(reserved)|\textcolor{red}{SPI\_USR\_CONF}}{1}{28}{{0}}%
\regfield{SPI\_SOFT\_RESET}{1}{27}{{0}}%
\regfield{SPI\_SLAVE\_MODE}{1}{26}{{0}}%
\regfield{(reserved)|\textcolor{red}{SPI\_DMA\_SEG\_MAGIC\_VALUE}}{4}{22}{{10}}%
\regfield{(reserved)}{10}{12}{0000000000}%
\regfield{SPI\_SLV\_WRBUF\_BITLEN\_EN}{1}{11}{{0}}%
\regfield{SPI\_SLV\_RDBUF\_BITLEN\_EN}{1}{10}{{0}}%
\regfield{SPI\_SLV\_WRDMA\_BITLEN\_EN}{1}{9}{{0}}%
\regfield{SPI\_SLV\_RDDMA\_BITLEN\_EN}{1}{8}{{0}}%
\regfield{(reserved)}{4}{4}{0000}%
\regfield{SPI\_RSCK\_DATA\_OUT}{1}{3}{{0}}%
\regfield{SPI\_CLK\_MODE\_13}{1}{2}{{0}}%
\regfield{SPI\_CLK\_MODE}{2}{0}{{0}}%
\reglabel{Reset}\regnewline%
\begin{regdesc}\begin{reglist}
\label{fielddesc:SPICLKMODE}\item [SPI\_CLK\_MODE] SPI 时钟模式控制位。可在 CONF 阶段配置。 (R/W)
\begin{itemize}
    \item 0：CS 信号无效时，SPI 时钟关闭；
    \item 1：CS 信号无效后，SPI 时钟延迟一个时钟周期；
    \item 2：CS 信号无效后，SPI 时钟延迟两个时钟周期；
    \item 3：SPI 时钟一直有效。
\end{itemize}
\label{fielddesc:SPICLKMODE13}\item [SPI\_CLK\_MODE\_13] 配置时钟模式。 (R/W)
\begin{itemize}
    \item 1：使用 SPI 时钟模式 1 或 3，在第一个跳变沿输出 B[0]/B[7] 数据；
    \item 0：使用 SPI 时钟模式 0 或 2，在第一个跳变沿输出 B[1]/B[6] 数据。
\end{itemize}
\label{fielddesc:SPIRSCKDATAOUT}\item [SPI\_RSCK\_DATA\_OUT] TSCK 与 RSCK 相同时，将节省半个周期。1：在 RSCK 上升沿输出数据。0：在 TSCK 上升沿输出数据。 (R/W)
\label{fielddesc:SPISLVRDDMABITLENEN}\item [SPI\_SLV\_RDDMA\_BITLEN\_EN] 置位此位，则在 DMA 控制的 Rd\_DMA 传输过程中，\hyperref[fielddesc:SPISLVDATABITLEN]{SPI\_SLV\_DATA\_BITLEN} 将用于存储 Rd\_DMA 传输的数据位长度。 (R/W)
\label{fielddesc:SPISLVWRDMABITLENEN}\item [SPI\_SLV\_WRDMA\_BITLEN\_EN] 置位此位，则在 DMA 控制的 Wr\_DMA 传输过程中，\hyperref[fielddesc:SPISLVDATABITLEN]{SPI\_SLV\_DATA\_BITLEN} 将用于存储 Wr\_DMA 传输的数据位长度。 (R/W)
\label{fielddesc:SPISLVRDBUFBITLENEN}\item [SPI\_SLV\_RDBUF\_BITLEN\_EN] 置位此位，则在 CPU 控制的 Rd\_BUF 传输过程中，\hyperref[fielddesc:SPISLVDATABITLEN]{SPI\_SLV\_DATA\_BITLEN} 将用于存储 Rd\_BUF 传输的数据位长度。 (R/W)
\label{fielddesc:SPISLVWRBUFBITLENEN}\item [SPI\_SLV\_WRBUF\_BITLEN\_EN] 置位此位，则在 CPU 控制的 Wr\_BUF 传输过程中，\hyperref[fielddesc:SPISLVDATABITLEN]{SPI\_SLV\_DATA\_BITLEN} 将用于存储 Wr\_BUF 传输的数据位长度。 (R/W)
\label{fielddesc:SPIDMASEGMAGICVALUE}\item [\textcolor{red}{SPI\_DMA\_SEG\_MAGIC\_VALUE（仅对 SPI2 有效）}] 在主机分段配置传输中，用于配置位图表的 Magic 值。 (R/W)
\label{fielddesc:SPISLAVEMODE}\item [SPI\_SLAVE\_MODE] 配置 SPI 工作模式。1：从机模式；0：主机模式。 (R/W)
\label{fielddesc:SPISOFTRESET}\item [SPI\_SOFT\_RESET] 软件置位使能位。置位此位，可复位 SPI 时钟线、CS 线和数据线。可在 CONF 阶段配置。 (WT)
\label{fielddesc:SPIUSRCONF}\item [\textcolor{red}{SPI\_USR\_CONF（仅对 SPI2 有效）}] 1：置位此位，使能当前分段配置传输的 CONF 阶段，开始分段配置传输。0：清除此位，则表明当前传输不是分段配置传输。 (R/W)
\end{reglist}\end{regdesc}
\end{register}


\begin{register}{H}{SPI\_SLAVE1\_REG}{0x{}00E4}\label{regdesc:SPISLAVE1REG}
\regfield{SPI\_SLV\_LAST\_ADDR}{6}{26}{{0}}%
\regfield{SPI\_SLV\_LAST\_COMMAND}{8}{18}{{0}}%
\regfield{SPI\_SLV\_DATA\_BITLEN}{18}{0}{{0}}%
\reglabel{Reset}\regnewline%
\begin{regdesc}\begin{reglist}
\label{fielddesc:SPISLVDATABITLEN}\item [SPI\_SLV\_DATA\_BITLEN] 在 SPI 从机全双工和半双工传输中，传输的数据位长。 (R/W/SS)
\label{fielddesc:SPISLVLASTCOMMAND}\item [SPI\_SLV\_LAST\_COMMAND] 从机模式下的命令值。 (R/W/SS)
\label{fielddesc:SPISLVLASTADDR}\item [SPI\_SLV\_LAST\_ADDR] 从机模式下的地址值。 (R/W/SS)
\end{reglist}\end{regdesc}
\end{register}


\begin{register}{H}{SPI\_CLOCK\_REG}{0x{}000C}\label{regdesc:SPICLOCKREG}
\regfield{SPI\_CLK\_EQU\_SYSCLK}{1}{31}{{1}}%
\regfield{(reserved)}{9}{22}{000000000}%
\regfield{SPI\_CLKDIV\_PRE}{4}{18}{{0}}%
\regfield{SPI\_CLKCNT\_N}{6}{12}{{0x{}3}}%
\regfield{SPI\_CLKCNT\_H}{6}{6}{{0x{}1}}%
\regfield{SPI\_CLKCNT\_L}{6}{0}{{0x{}3}}%
\reglabel{Reset}\regnewline%
\begin{regdesc}\begin{reglist}
\label{fielddesc:SPICLKCNTL}\item [SPI\_CLKCNT\_L] 主机模式下，必须与 SPI\_CLKCNT\_N 相等。在从机模式下，必须为 0。可在 CONF 阶段配置。 (R/W)
\label{fielddesc:SPICLKCNTH}\item [SPI\_CLKCNT\_H] 主机模式下，必须为 (SPI\_CLKCNT\_N + 1)/2 - 1 的向下取整值。在从机模式下，必须为 0。可在 CONF 阶段配置。 (R/W)
\label{fielddesc:SPICLKCNTN}\item [SPI\_CLKCNT\_N] 主机模式下，SPI\_CLK 的分频系数。因此 SPI\_CLK 频率为 f{\textrm{apb\_clk}}/(SPI\_CLKDIV\_PRE + 1)/(SPI\_CLKCNT\_N + 1)。可在 CONF 阶段配置。 (R/W)
\label{fielddesc:SPICLKDIVPRE}\item [SPI\_CLKDIV\_PRE] 主机模式下，SPI\_CLK 的预分频系数。可在 CONF 阶段配置。 (R/W)
\label{fielddesc:SPICLKEQUSYSCLK}\item [SPI\_CLK\_EQU\_SYSCLK] 在主机模式下，1：SPI\_CLK 与 APB\_CLK 频率相同；0：SPI\_CLK 为 APB\_CLK 的分频时钟。可在 CONF 阶段配置。 (R/W)
\end{reglist}\end{regdesc}
\end{register}


\begin{register}{H}{SPI\_CLK\_GATE\_REG}{0x{}00E8}\label{regdesc:SPICLKGATEREG}
\regfield{(reserved)}{29}{3}{00000000000000000000000000000}%
\regfield{SPI\_MST\_CLK\_SEL}{1}{2}{{0}}%
\regfield{SPI\_MST\_CLK\_ACTIVE}{1}{1}{{0}}%
\regfield{SPI\_CLK\_EN}{1}{0}{{0}}%
\reglabel{Reset}\regnewline%
\begin{regdesc}\begin{reglist}
\label{fielddesc:SPICLKEN}\item [SPI\_CLK\_EN] 置位此位，使能时钟门控。 (R/W)
\label{fielddesc:SPIMSTCLKACTIVE}\item [SPI\_MST\_CLK\_ACTIVE] 置位此位，SPI 模块时钟上电。 (R/W)
\label{fielddesc:SPIMSTCLKSEL}\item [SPI\_MST\_CLK\_SEL] 该位用于选择主机模式下 SPI 模块的时钟源。1：选择 PLL\_F80M\_CLK；0：选择 XTAL\_CLK。 (R/W)
\end{reglist}\end{regdesc}
\end{register}


\begin{register}{H}{\textcolor{red}{SPI\_DIN\_MODE\_REG}}{0x{}0024}\label{regdesc:SPIDINMODEREG}
\regfield{(reserved)}{15}{17}{000000000000000}%
\regfield{SPI\_TIMING\_HCLK\_ACTIVE}{1}{16}{{0}}%
\regfield{(reserved)|\textcolor{red}{SPI\_DIN7\_MODE}}{2}{14}{{0}}%
\regfield{(reserved)|\textcolor{red}{SPI\_DIN6\_MODE}}{2}{12}{{0}}%
\regfield{(reserved)|\textcolor{red}{SPI\_DIN5\_MODE}}{2}{10}{{0}}%
\regfield{(reserved)|\textcolor{red}{SPI\_DIN4\_MODE}}{2}{8}{{0}}%
\regfield{SPI\_DIN3\_MODE}{2}{6}{{0}}%
\regfield{SPI\_DIN2\_MODE}{2}{4}{{0}}%
\regfield{SPI\_DIN1\_MODE}{2}{2}{{0}}%
\regfield{SPI\_DIN0\_MODE}{2}{0}{{0}}%
\reglabel{Reset}\regnewline%
\begin{regdesc}\begin{reglist}
\label{fielddesc:SPIDIN0MODE}\item [SPI\_DIN0\_MODE] 配置输入数据 bit0 的输入模式。可在 CONF 阶段配置。 (R/W)
\begin{itemize}
    \item 0：无输入延迟
    \item 1：输入数据在 \hyperref[clk:SPICLK]{SPI\_CLK} 下降沿延迟 (\hyperref[fielddesc:SPIDIN0NUM]{SPI\_DIN0\_NUM} +1) 个周期
    \item 2：输入数据在 \hyperref[clk:SPICLK]{SPI\_CLK} 上升沿延迟 (\hyperref[fielddesc:SPIDIN0NUM]{SPI\_DIN0\_NUM} + 1) 个周期
    \item 3：输入数据在 \hyperref[clk:clkhclk]{clk\_hclk} 上升沿延迟 (\hyperref[fielddesc:SPIDIN0NUM]{SPI\_DIN0\_NUM} + 1) 个周期，然后再在 \hyperref[clk:SPICLK]{SPI\_CLK} 下降沿延迟一个时钟周期
\end{itemize}

\label{fielddesc:SPIDIN1MODE}\item [SPI\_DIN1\_MODE] 配置输入数据 bit1 的输入模式。可在 CONF 阶段配置。 (R/W)
\begin{itemize}
    \item 0：无输入延迟
    \item 1：输入数据在 \hyperref[clk:SPICLK]{SPI\_CLK} 下降沿延迟 (\hyperref[fielddesc:SPIDIN0NUM]{SPI\_DIN1\_NUM} +1) 个周期
    \item 2：输入数据在 \hyperref[clk:SPICLK]{SPI\_CLK} 上升沿延迟 (\hyperref[fielddesc:SPIDIN0NUM]{SPI\_DIN1\_NUM} + 1) 个周期
    \item 3：输入数据在 \hyperref[clk:clkhclk]{clk\_hclk} 上升沿延迟 (\hyperref[fielddesc:SPIDIN0NUM]{SPI\_DIN1\_NUM} + 1) 个周期，然后再在 \hyperref[clk:SPICLK]{SPI\_CLK} 下降沿延迟一个时钟周期
\end{itemize}
\label{fielddesc:SPIDIN2MODE}\item [SPI\_DIN2\_MODE] 配置输入数据 bit2 的输入模式。可在 CONF 阶段配置。 (R/W)
\begin{itemize}
    \item 0：无输入延迟
    \item 1：输入数据在 \hyperref[clk:SPICLK]{SPI\_CLK} 下降沿延迟 (\hyperref[fielddesc:SPIDIN0NUM]{SPI\_DIN2\_NUM} +1) 个周期
    \item 2：输入数据在 \hyperref[clk:SPICLK]{SPI\_CLK} 上升沿延迟 (\hyperref[fielddesc:SPIDIN0NUM]{SPI\_DIN2\_NUM} + 1) 个周期
    \item 3：输入数据在 \hyperref[clk:clkhclk]{clk\_hclk} 上升沿延迟 (\hyperref[fielddesc:SPIDIN0NUM]{SPI\_DIN2\_NUM} + 1) 个周期，然后再在 \hyperref[clk:SPICLK]{SPI\_CLK} 下降沿延迟一个时钟周期
\end{itemize}
\label{fielddesc:SPIDIN3MODE}\item [SPI\_DIN3\_MODE] 配置输入数据 bit3 的输入模式。可在 CONF 阶段配置。 (R/W)
\begin{itemize}
    \item 0：无输入延迟
    \item 1：输入数据在 \hyperref[clk:SPICLK]{SPI\_CLK} 下降沿延迟 (\hyperref[fielddesc:SPIDIN0NUM]{SPI\_DIN3\_NUM} +1) 个周期
    \item 2：输入数据在 \hyperref[clk:SPICLK]{SPI\_CLK} 上升沿延迟 (\hyperref[fielddesc:SPIDIN0NUM]{SPI\_DIN3\_NUM} + 1) 个周期
    \item 3：输入数据在 \hyperref[clk:clkhclk]{clk\_hclk} 上升沿延迟 (\hyperref[fielddesc:SPIDIN0NUM]{SPI\_DIN3\_NUM} + 1) 个周期，然后再在 \hyperref[clk:SPICLK]{SPI\_CLK} 下降沿延迟一个时钟周期
\end{itemize}

\item [见下页]
\end{reglist}\end{regdesc}
\end{register}

\addtocounter{Regfloat}{-1}
\begin{register}{H}{\textcolor{red}{SPI\_DIN\_MODE\_REG}}{0x{}0024}
\begin{regdesc}\begin{reglist}
\item [接上页]

\label{fielddesc:SPIDIN4MODE}\item [\textcolor{red}{SPI\_DIN4\_MODE（仅对 SPI2 有效）}] 配置输入数据 bit4 的输入模式。可在 CONF 阶段配置。 (R/W)
\begin{itemize}
    \item 0：无输入延迟
    \item 1：输入数据在 \hyperref[clk:SPICLK]{SPI\_CLK} 下降沿延迟 (\hyperref[fielddesc:SPIDIN0NUM]{SPI\_DIN4\_NUM} +1) 个周期
    \item 2：输入数据在 \hyperref[clk:SPICLK]{SPI\_CLK} 上升沿延迟 (\hyperref[fielddesc:SPIDIN0NUM]{SPI\_DIN4\_NUM} + 1) 个周期
    \item 3：输入数据在 \hyperref[clk:clkhclk]{clk\_hclk} 上升沿延迟 (\hyperref[fielddesc:SPIDIN0NUM]{SPI\_DIN4\_NUM} + 1) 个周期，然后再在 \hyperref[clk:SPICLK]{SPI\_CLK} 下降沿延迟一个时钟周期
\end{itemize}
\label{fielddesc:SPIDIN5MODE}\item [\textcolor{red}{SPI\_DIN5\_MODE（仅对 SPI2 有效）}] 配置输入数据 bit5 的输入模式。可在 CONF 阶段配置。 (R/W)
\begin{itemize}
    \item 0：无输入延迟
    \item 1：输入数据在 \hyperref[clk:SPICLK]{SPI\_CLK} 下降沿延迟 (\hyperref[fielddesc:SPIDIN0NUM]{SPI\_DIN5\_NUM} +1) 个周期
    \item 2：输入数据在 \hyperref[clk:SPICLK]{SPI\_CLK} 上升沿延迟 (\hyperref[fielddesc:SPIDIN0NUM]{SPI\_DIN5\_NUM} + 1) 个周期
    \item 3：输入数据在 \hyperref[clk:clkhclk]{clk\_hclk} 上升沿延迟 (\hyperref[fielddesc:SPIDIN0NUM]{SPI\_DIN5\_NUM} + 1) 个周期，然后再在 \hyperref[clk:SPICLK]{SPI\_CLK} 下降沿延迟一个时钟周期
\end{itemize}
\label{fielddesc:SPIDIN6MODE}\item [\textcolor{red}{SPI\_DIN6\_MODE（仅对 SPI2 有效）}] 配置输入数据 bit6 的输入模式。可在 CONF 阶段配置。 (R/W)
\begin{itemize}
    \item 0：无输入延迟
    \item 1：输入数据在 \hyperref[clk:SPICLK]{SPI\_CLK} 下降沿延迟 (\hyperref[fielddesc:SPIDIN0NUM]{SPI\_DIN6\_NUM} +1) 个周期
    \item 2：输入数据在 \hyperref[clk:SPICLK]{SPI\_CLK} 上升沿延迟 (\hyperref[fielddesc:SPIDIN0NUM]{SPI\_DIN6\_NUM} + 1) 个周期
    \item 3：输入数据在 \hyperref[clk:clkhclk]{clk\_hclk} 上升沿延迟 (\hyperref[fielddesc:SPIDIN0NUM]{SPI\_DIN6\_NUM} + 1) 个周期，然后再在 \hyperref[clk:SPICLK]{SPI\_CLK} 下降沿延迟一个时钟周期
\end{itemize}
\label{fielddesc:SPIDIN7MODE}\item [\textcolor{red}{SPI\_DIN7\_MODE（仅对 SPI2 有效）}] 配置输入数据 bit7 的输入模式。可在 CONF 阶段配置。 (R/W)
\begin{itemize}
    \item 0：无输入延迟
    \item 1：输入数据在 \hyperref[clk:SPICLK]{SPI\_CLK} 下降沿延迟 (\hyperref[fielddesc:SPIDIN0NUM]{SPI\_DIN7\_NUM} +1) 个周期
    \item 2：输入数据在 \hyperref[clk:SPICLK]{SPI\_CLK} 上升沿延迟 (\hyperref[fielddesc:SPIDIN0NUM]{SPI\_DIN7\_NUM} + 1) 个周期
    \item 3：输入数据在 \hyperref[clk:clkhclk]{clk\_hclk} 上升沿延迟 (\hyperref[fielddesc:SPIDIN0NUM]{SPI\_DIN7\_NUM} + 1) 个周期，然后再在 \hyperref[clk:SPICLK]{SPI\_CLK} 下降沿延迟一个时钟周期
\end{itemize}
\label{fielddesc:SPITIMINGHCLKACTIVE}\item [SPI\_TIMING\_HCLK\_ACTIVE] 1：使能 SPI 输入信号时序模块的高频时钟 clk\_hclk；0：禁用 clk\_hclk。可在 CONF 阶段配置。 (R/W)
\end{reglist}\end{regdesc}
\end{register}


\begin{register}{H}{\textcolor{red}{SPI\_DIN\_NUM\_REG}}{0x{}0028}\label{regdesc:SPIDINNUMREG}
\regfield{(reserved)}{16}{16}{0000000000000000}%
\regfield{(reserved)|\textcolor{red}{SPI\_DIN7\_NUM}}{2}{14}{{0}}%
\regfield{(reserved)|\textcolor{red}{SPI\_DIN6\_NUM}}{2}{12}{{0}}%
\regfield{(reserved)|\textcolor{red}{SPI\_DIN5\_NUM}}{2}{10}{{0}}%
\regfield{(reserved)|\textcolor{red}{SPI\_DIN4\_NUM}}{2}{8}{{0}}%
\regfield{SPI\_DIN3\_NUM}{2}{6}{{0}}%
\regfield{SPI\_DIN2\_NUM}{2}{4}{{0}}%
\regfield{SPI\_DIN1\_NUM}{2}{2}{{0}}%
\regfield{SPI\_DIN0\_NUM}{2}{0}{{0}}%
\reglabel{Reset}\regnewline%
\begin{regdesc}\begin{reglist}
\label{fielddesc:SPIDIN0NUM}\item [SPI\_DIN0\_NUM] 配置输入信号 bit0 的延迟周期数。具体的延迟模式由 \hyperref[fielddesc:SPIDIN0MODE]{SPI\_DIN0\_MODE} 配置。可在 CONF 阶段配置。 (R/W)
\label{fielddesc:SPIDIN1NUM}\item [SPI\_DIN1\_NUM] 配置输入信号 bit1 的延迟周期数。具体的延迟模式由 \hyperref[fielddesc:SPIDIN2MODE]{SPI\_DIN2\_MODE} 配置。可在 CONF 阶段配置。 (R/W)
\label{fielddesc:SPIDIN2NUM}\item [SPI\_DIN2\_NUM] 配置输入信号 bit2 的延迟周期数。具体的延迟模式由 \hyperref[fielddesc:SPIDIN0MODE]{SPI\_DIN0\_MODE} 配置。可在 CONF 阶段配置。 (R/W)
\label{fielddesc:SPIDIN3NUM}\item [SPI\_DIN3\_NUM] 配置输入信号 bit3 的延迟周期数。具体的延迟模式由 \hyperref[fielddesc:SPIDIN3MODE]{SPI\_DIN3\_MODE} 配置。可在 CONF 阶段配置。 (R/W)
\label{fielddesc:SPIDIN4NUM}\item [\textcolor{red}{SPI\_DIN4\_NUM（仅对 SPI2 有效）}] 配置输入信号 bit4 的延迟周期数。具体的延迟模式由 \hyperref[fielddesc:SPIDIN4MODE]{SPI\_DIN4\_MODE} 配置。可在 CONF 阶段配置。 (R/W)
\label{fielddesc:SPIDIN5NUM}\item [\textcolor{red}{SPI\_DIN5\_NUM（仅对 SPI2 有效）}] 配置输入信号 bit5 的延迟周期数。具体的延迟模式由 \hyperref[fielddesc:SPIDIN5MODE]{SPI\_DIN5\_MODE} 配置。可在 CONF 阶段配置。 (R/W)
\label{fielddesc:SPIDIN6NUM}\item [\textcolor{red}{SPI\_DIN6\_NUM（仅对 SPI2 有效）}] 配置输入信号 bit6 的延迟周期数。具体的延迟模式由 \hyperref[fielddesc:SPIDIN6MODE]{SPI\_DIN6\_MODE} 配置。可在 CONF 阶段配置。 (R/W)
\label{fielddesc:SPIDIN7NUM}\item [\textcolor{red}{SPI\_DIN7\_NUM（仅对 SPI2 有效）}] 配置输入信号 bit7 的延迟周期数。具体的延迟模式由 \hyperref[fielddesc:SPIDIN7MODE]{SPI\_DIN7\_MODE} 配置。可在 CONF 阶段配置。 (R/W)
\end{reglist}\end{regdesc}
\end{register}


\begin{register}{H}{\textcolor{red}{SPI\_DOUT\_MODE\_REG}}{0x{}002C}\label{regdesc:SPIDOUTMODEREG}
\regfield{(reserved)}{23}{9}{00000000000000000000000}%
\regfield{(reserved)|\textcolor{red}{SPI\_D\_DQS\_MODE}}{1}{8}{{0}}%
\regfield{(reserved)|\textcolor{red}{SPI\_DOUT7\_MODE}}{1}{7}{{0}}%
\regfield{(reserved)|\textcolor{red}{SPI\_DOUT6\_MODE}}{1}{6}{{0}}%
\regfield{(reserved)|\textcolor{red}{SPI\_DOUT5\_MODE}}{1}{5}{{0}}%
\regfield{(reserved)|\textcolor{red}{SPI\_DOUT4\_MODE}}{1}{4}{{0}}%
\regfield{SPI\_DOUT3\_MODE}{1}{3}{{0}}%
\regfield{SPI\_DOUT2\_MODE}{1}{2}{{0}}%
\regfield{SPI\_DOUT1\_MODE}{1}{1}{{0}}%
\regfield{SPI\_DOUT0\_MODE}{1}{0}{{0}}%
\reglabel{Reset}\regnewline%
\begin{regdesc}\begin{reglist}
\label{fielddesc:SPIDOUT0MODE}\item [SPI\_DOUT0\_MODE] 配置输出数据 bit0 信号的输出模式。可在 CONF 阶段配置。 (R/W)
\begin{itemize}
    \item 0：无输出延迟
    \item 1：在 \hyperref[clk:SPICLK]{SPI\_CLK} 时钟下降沿，延迟一个时钟周期后输出
\end{itemize}
\label{fielddesc:SPIDOUT1MODE}\item [SPI\_DOUT1\_MODE] 配置输出数据 bit1 信号的输出模式。可在 CONF 阶段配置。 (R/W)
\begin{itemize}
    \item 0：无输出延迟
    \item 1：在 \hyperref[clk:SPICLK]{SPI\_CLK} 时钟下降沿，延迟一个时钟周期后输出
\end{itemize}
\label{fielddesc:SPIDOUT2MODE}\item [SPI\_DOUT2\_MODE] 配置输出数据 bit2 信号的输出模式。可在 CONF 阶段配置。 (R/W)
\begin{itemize}
    \item 0：无输出延迟
    \item 1：在 \hyperref[clk:SPICLK]{SPI\_CLK} 时钟下降沿，延迟一个时钟周期后输出
\end{itemize}
\label{fielddesc:SPIDOUT3MODE}\item [SPI\_DOUT3\_MODE] 配置输出数据 bit3 信号的输出模式。可在 CONF 阶段配置。 (R/W)
\begin{itemize}
    \item 0：无输出延迟
    \item 1：在 \hyperref[clk:SPICLK]{SPI\_CLK} 时钟下降沿，延迟一个时钟周期后输出
\end{itemize}
\label{fielddesc:SPIDOUT4MODE}\item [\textcolor{red}{SPI\_DOUT4\_MODE（仅对 SPI2 有效）}] 配置输出数据 bit4 信号的输出模式。可在 CONF 阶段配置。 (R/W)
\begin{itemize}
    \item 0：无输出延迟
    \item 1：在 \hyperref[clk:SPICLK]{SPI\_CLK} 时钟下降沿，延迟一个时钟周期后输出
\end{itemize}
\label{fielddesc:SPIDOUT5MODE}\item [\textcolor{red}{SPI\_DOUT5\_MODE（仅对 SPI2 有效）}] 配置输出数据 bit5 信号的输出模式。可在 CONF 阶段配置。 (R/W)
\begin{itemize}
    \item 0：无输出延迟
    \item 1：在 \hyperref[clk:SPICLK]{SPI\_CLK} 时钟下降沿，延迟一个时钟周期后输出
\end{itemize}
\label{fielddesc:SPIDOUT6MODE}\item [\textcolor{red}{SPI\_DOUT6\_MODE（仅对 SPI2 有效）}] 配置输出数据 bit6 信号的输出模式。可在 CONF 阶段配置。 (R/W)
\begin{itemize}
    \item 0：无输出延迟
    \item 1：在 \hyperref[clk:SPICLK]{SPI\_CLK} 时钟下降沿，延迟一个时钟周期后输出
\end{itemize}

\item [见下页]
\end{reglist}\end{regdesc}
\end{register}

\addtocounter{Regfloat}{-1}
\begin{register}{H}{\textcolor{red}{SPI\_DOUT\_MODE\_REG}}{0x{}002C}
\begin{regdesc}\begin{reglist}
\item [接上页]

\label{fielddesc:SPIDOUT7MODE}\item [\textcolor{red}{SPI\_DOUT7\_MODE（仅对 SPI2 有效）}] 配置输出数据 bit7 信号的输出模式。可在 CONF 阶段配置。 (R/W)
\begin{itemize}
    \item 0：无输出延迟
    \item 1：在 \hyperref[clk:SPICLK]{SPI\_CLK} 时钟下降沿，延迟一个时钟周期后输出
\end{itemize}
\label{fielddesc:SPIDDQSMODE}\item [\textcolor{red}{SPI\_D\_DQS\_MODE（仅对 SPI2 有效）}] 配置输出信号 SPI\_DQS 的输出模式。可在 CONF 阶段配置。 (R/W)
\begin{itemize}
    \item 0：无输出延迟
    \item 1：在 \hyperref[clk:SPICLK]{SPI\_CLK} 时钟下降沿，延迟一个时钟周期后输出
\end{itemize}
\end{reglist}\end{regdesc}
\end{register}


\begin{register}{H}{\textcolor{red}{SPI\_DMA\_INT\_ENA\_REG}}{0x{}0034}\label{regdesc:SPIDMAINTENAREG}
\regfield{(reserved)}{11}{21}{00000000000}%
\regfield{SPI\_APP1\_INT\_ENA}{1}{20}{{0}}%
\regfield{SPI\_APP2\_INT\_ENA}{1}{19}{{0}}%
\regfield{SPI\_MST\_TX\_AFIFO\_REMPTY\_ERR\_INT\_ENA}{1}{18}{{0}}%
\regfield{SPI\_MST\_RX\_AFIFO\_WFULL\_ERR\_INT\_ENA}{1}{17}{{0}}%
\regfield{SPI\_SLV\_CMD\_ERR\_INT\_ENA}{1}{16}{{0}}%
\regfield{(reserved)}{1}{15}{{0}}%
\regfield{(reserved)|\textcolor{red}{SPI\_SEG\_MAGIC\_ERR\_INT\_ENA}}{1}{14}{{0}}%
\regfield{SPI\_DMA\_SEG\_TRANS\_DONE\_INT\_ENA}{1}{13}{{0}}%
\regfield{SPI\_TRANS\_DONE\_INT\_ENA}{1}{12}{{0}}%
\regfield{SPI\_SLV\_WR\_BUF\_DONE\_INT\_ENA}{1}{11}{{0}}%
\regfield{SPI\_SLV\_RD\_BUF\_DONE\_INT\_ENA}{1}{10}{{0}}%
\regfield{SPI\_SLV\_WR\_DMA\_DONE\_INT\_ENA}{1}{9}{{0}}%
\regfield{SPI\_SLV\_RD\_DMA\_DONE\_INT\_ENA}{1}{8}{{0}}%
\regfield{SPI\_SLV\_CMDA\_INT\_ENA}{1}{7}{{0}}%
\regfield{SPI\_SLV\_CMD9\_INT\_ENA}{1}{6}{{0}}%
\regfield{SPI\_SLV\_CMD8\_INT\_ENA}{1}{5}{{0}}%
\regfield{SPI\_SLV\_CMD7\_INT\_ENA}{1}{4}{{0}}%
\regfield{SPI\_SLV\_EN\_QPI\_INT\_ENA}{1}{3}{{0}}%
\regfield{SPI\_SLV\_EX\_QPI\_INT\_ENA}{1}{2}{{0}}%
\regfield{SPI\_DMA\_OUTFIFO\_EMPTY\_ERR\_INT\_ENA}{1}{1}{{0}}%
\regfield{SPI\_DMA\_INFIFO\_FULL\_ERR\_INT\_ENA}{1}{0}{{0}}%
\reglabel{Reset}\regnewline%
\begin{regdesc}\begin{reglist}
\label{fielddesc:SPIDMAINFIFOFULLERRINTENA}\item [SPI\_DMA\_INFIFO\_FULL\_ERR\_INT\_ENA] \hyperref[int:SPIINFIFOFULLERRINT]{SPI\_DMA\_INFIFO\_FULL\_ERR\_INT} 的中断使能位。 (R/W)
\label{fielddesc:SPIDMAOUTFIFOEMPTYERRINTENA}\item [SPI\_DMA\_OUTFIFO\_EMPTY\_ERR\_INT\_ENA] \hyperref[int:SPIOUTFIFOEMPTYERRINT]{SPI\_DMA\_OUTFIFO\_EMPTY\_ERR\_INT} 的中断使能位。 (R/W)
\label{fielddesc:SPISLVEXQPIINTENA}\item [SPI\_SLV\_EX\_QPI\_INT\_ENA] \hyperref[int:SPISLVEXQPIINT]{SPI\_SLV\_EX\_QPI\_INT} 的中断使能位。 (R/W)
\label{fielddesc:SPISLVENQPIINTENA}\item [SPI\_SLV\_EN\_QPI\_INT\_ENA] \hyperref[int:SPISLVENQPIINT]{SPI\_SLV\_EN\_QPI\_INT} 的中断使能位。 (R/W)
\label{fielddesc:SPISLVCMD7INTENA}\item [SPI\_SLV\_CMD7\_INT\_ENA] \hyperref[int:SPISLVCMD7INT]{SPI\_SLV\_CMD7\_INT} 的中断使能位。 (R/W)
\label{fielddesc:SPISLVCMD8INTENA}\item [SPI\_SLV\_CMD8\_INT\_ENA] \hyperref[int:SPISLVCMD8INT]{SPI\_SLV\_CMD8\_INT} 的中断使能位。 (R/W)
\label{fielddesc:SPISLVCMD9INTENA}\item [SPI\_SLV\_CMD9\_INT\_ENA] \hyperref[int:SPISLVCMD9INT]{SPI\_SLV\_CMD9\_INT} 的中断使能位。 (R/W)
\label{fielddesc:SPISLVCMDAINTENA}\item [SPI\_SLV\_CMDA\_INT\_ENA] \hyperref[int:SPISLVCMDAINT]{SPI\_SLV\_CMDA\_INT} 的中断使能位。 (R/W)
\label{fielddesc:SPISLVRDDMADONEINTENA}\item [SPI\_SLV\_RD\_DMA\_DONE\_INT\_ENA] \hyperref[int:SPISLVRDDMADONEINT]{SPI\_SLV\_RD\_DMA\_DONE\_INT} 的中断使能位。 (R/W)
\label{fielddesc:SPISLVWRDMADONEINTENA}\item [SPI\_SLV\_WR\_DMA\_DONE\_INT\_ENA] \hyperref[int:SPISLVWRDMADONEINT]{SPI\_SLV\_WR\_DMA\_DONE\_INT} 的中断使能位。 (R/W)
\label{fielddesc:SPISLVRDBUFDONEINTENA}\item [SPI\_SLV\_RD\_BUF\_DONE\_INT\_ENA] \hyperref[int:SPISLVRDBUFDONEINT]{SPI\_SLV\_RD\_BUF\_DONE\_INT} 的中断使能位。 (R/W)
\label{fielddesc:SPISLVWRBUFDONEINTENA}\item [SPI\_SLV\_WR\_BUF\_DONE\_INT\_ENA] \hyperref[int:SPISLVWRBUFDONEINT]{SPI\_SLV\_WR\_BUF\_DONE\_INT} 的中断使能位。 (R/W)
\label{fielddesc:SPITRANSDONEINTENA}\item [SPI\_TRANS\_DONE\_INT\_ENA] \hyperref[int:SPITRANSDONEINT]{SPI\_TRANS\_DONE\_INT} 的中断使能位。 (R/W)
\label{fielddesc:SPIDMASEGTRANSDONEINTENA}\item [SPI\_DMA\_SEG\_TRANS\_DONE\_INT\_ENA] \hyperref[int:SPIDMASEGTRANSDONEINT]{SPI\_DMA\_SEG\_TRANS\_DONE\_INT} 的中断使能位。 (R/W)
\label{fielddesc:SPISEGMAGICERRINTENA}\item [\textcolor{red}{SPI\_SEG\_MAGIC\_ERR\_INT\_ENA（仅对 SPI2 有效）}] \hyperref[int:SPISEGMAGICERRINT]{SPI\_SEG\_MAGIC\_ERR\_INT} 的中断使能位。 (R/W)
\label{fielddesc:SPISLVCMDERRINTENA}\item [SPI\_SLV\_CMD\_ERR\_INT\_ENA] \hyperref[int:SPISLVCMDERRINT]{SPI\_SLV\_CMD\_ERR\_INT} 的中断使能位。 (R/W)

\item [见下页]
\end{reglist}\end{regdesc}
\end{register}

\addtocounter{Regfloat}{-1}
\begin{register}{H}{\textcolor{red}{SPI\_DMA\_INT\_ENA\_REG}}{0x{}0034}
\begin{regdesc}\begin{reglist}
\item [接上页]
\label{fielddesc:SPIMSTRXAFIFOWFULLERRINTENA}\item [SPI\_MST\_RX\_AFIFO\_WFULL\_ERR\_INT\_ENA] \hyperref[int:SPIMSTRXAFIFOWFULLERRINT]{SPI\_MST\_RX\_AFIFO\_WFULL\_ERR\_INT} 的中断使能位。 (R/W)
\label{fielddesc:SPIMSTTXAFIFOREMPTYERRINTENA}\item [SPI\_MST\_TX\_AFIFO\_REMPTY\_ERR\_INT\_ENA] \hyperref[int:SPIMSTTXAFIFOREMPTYERRINT]{SPI\_MST\_TX\_AFIFO\_REMPTY\_ERR\_INT} 的中断使能位。 (R/W)
\label{fielddesc:SPIAPP2INTENA}\item [SPI\_APP2\_INT\_ENA] \hyperref[int:SPIAPP2INT]{SPI\_APP2\_INT} 的中断使能位。 (R/W)
\label{fielddesc:SPIAPP1INTENA}\item [SPI\_APP1\_INT\_ENA] \hyperref[int:SPIAPP1INT]{SPI\_APP1\_INT} 的中断使能位。 (R/W)
\end{reglist}\end{regdesc}
\end{register}


\begin{register}{H}{\textcolor{red}{SPI\_DMA\_INT\_CLR\_REG}}{0x{}0038}\label{regdesc:SPIDMAINTCLRREG}
\regfield{(reserved)}{11}{21}{00000000000}%
\regfield{SPI\_APP1\_INT\_CLR}{1}{20}{{0}}%
\regfield{SPI\_APP2\_INT\_CLR}{1}{19}{{0}}%
\regfield{SPI\_MST\_TX\_AFIFO\_REMPTY\_ERR\_INT\_CLR}{1}{18}{{0}}%
\regfield{SPI\_MST\_RX\_AFIFO\_WFULL\_ERR\_INT\_CLR}{1}{17}{{0}}%
\regfield{SPI\_SLV\_CMD\_ERR\_INT\_CLR}{1}{16}{{0}}%
\regfield{(reserved)}{1}{15}{{0}}%
\regfield{(reserved)|\textcolor{red}{SPI\_SEG\_MAGIC\_ERR\_INT\_CLR}}{1}{14}{{0}}%
\regfield{SPI\_DMA\_SEG\_TRANS\_DONE\_INT\_CLR}{1}{13}{{0}}%
\regfield{SPI\_TRANS\_DONE\_INT\_CLR}{1}{12}{{0}}%
\regfield{SPI\_SLV\_WR\_BUF\_DONE\_INT\_CLR}{1}{11}{{0}}%
\regfield{SPI\_SLV\_RD\_BUF\_DONE\_INT\_CLR}{1}{10}{{0}}%
\regfield{SPI\_SLV\_WR\_DMA\_DONE\_INT\_CLR}{1}{9}{{0}}%
\regfield{SPI\_SLV\_RD\_DMA\_DONE\_INT\_CLR}{1}{8}{{0}}%
\regfield{SPI\_SLV\_CMDA\_INT\_CLR}{1}{7}{{0}}%
\regfield{SPI\_SLV\_CMD9\_INT\_CLR}{1}{6}{{0}}%
\regfield{SPI\_SLV\_CMD8\_INT\_CLR}{1}{5}{{0}}%
\regfield{SPI\_SLV\_CMD7\_INT\_CLR}{1}{4}{{0}}%
\regfield{SPI\_SLV\_EN\_QPI\_INT\_CLR}{1}{3}{{0}}%
\regfield{SPI\_SLV\_EX\_QPI\_INT\_CLR}{1}{2}{{0}}%
\regfield{SPI\_DMA\_OUTFIFO\_EMPTY\_ERR\_INT\_CLR}{1}{1}{{0}}%
\regfield{SPI\_DMA\_INFIFO\_FULL\_ERR\_INT\_CLR}{1}{0}{{0}}%
\reglabel{Reset}\regnewline%
\begin{regdesc}\begin{reglist}
\label{fielddesc:SPIDMAINFIFOFULLERRINTCLR}\item [SPI\_DMA\_INFIFO\_FULL\_ERR\_INT\_CLR] \hyperref[int:SPIINFIFOFULLERRINT]{SPI\_DMA\_INFIFO\_FULL\_ERR\_INT} 的中断清除位。 (WT)
\label{fielddesc:SPIDMAOUTFIFOEMPTYERRINTCLR}\item [SPI\_DMA\_OUTFIFO\_EMPTY\_ERR\_INT\_CLR] \hyperref[int:SPIOUTFIFOEMPTYERRINT]{SPI\_DMA\_OUTFIFO\_EMPTY\_ERR\_INT} 的中断清除位。 (WT)
\label{fielddesc:SPISLVEXQPIINTCLR}\item [SPI\_SLV\_EX\_QPI\_INT\_CLR] \hyperref[int:SPISLVEXQPIINT]{SPI\_SLV\_EX\_QPI\_INT} 的中断清除位。 (WT)
\label{fielddesc:SPISLVENQPIINTCLR}\item [SPI\_SLV\_EN\_QPI\_INT\_CLR] \hyperref[int:SPISLVENQPIINT]{SPI\_SLV\_EN\_QPI\_INT} 的中断清除位。 (WT)
\label{fielddesc:SPISLVCMD7INTCLR}\item [SPI\_SLV\_CMD7\_INT\_CLR] \hyperref[int:SPISLVCMD7INT]{SPI\_SLV\_CMD7\_INT} 的中断清除位。 (WT)
\label{fielddesc:SPISLVCMD8INTCLR}\item [SPI\_SLV\_CMD8\_INT\_CLR] \hyperref[int:SPISLVCMD8INT]{SPI\_SLV\_CMD8\_INT} 的中断清除位。 (WT)
\label{fielddesc:SPISLVCMD9INTCLR}\item [SPI\_SLV\_CMD9\_INT\_CLR] \hyperref[int:SPISLVCMD9INT]{SPI\_SLV\_CMD9\_INT} 的中断清除位。 (WT)
\label{fielddesc:SPISLVCMDAINTCLR}\item [SPI\_SLV\_CMDA\_INT\_CLR] \hyperref[int:SPISLVCMDAINT]{SPI\_SLV\_CMDA\_INT} 的中断清除位。 (WT)
\label{fielddesc:SPISLVRDDMADONEINTCLR}\item [SPI\_SLV\_RD\_DMA\_DONE\_INT\_CLR] \hyperref[int:SPISLVRDDMADONEINT]{SPI\_SLV\_RD\_DMA\_DONE\_INT} 的中断清除位。 (WT)
\label{fielddesc:SPISLVWRDMADONEINTCLR}\item [SPI\_SLV\_WR\_DMA\_DONE\_INT\_CLR] \hyperref[int:SPISLVWRDMADONEINT]{SPI\_SLV\_WR\_DMA\_DONE\_INT} 的中断清除位。 (WT)
\label{fielddesc:SPISLVRDBUFDONEINTCLR}\item [SPI\_SLV\_RD\_BUF\_DONE\_INT\_CLR] \hyperref[int:SPISLVRDBUFDONEINT]{SPI\_SLV\_RD\_BUF\_DONE\_INT} 的中断清除位。 (WT)
\label{fielddesc:SPISLVWRBUFDONEINTCLR}\item [SPI\_SLV\_WR\_BUF\_DONE\_INT\_CLR] \hyperref[int:SPISLVWRBUFDONEINT]{SPI\_SLV\_WR\_BUF\_DONE\_INT} 的中断清除位。 (WT)
\label{fielddesc:SPITRANSDONEINTCLR}\item [SPI\_TRANS\_DONE\_INT\_CLR] \hyperref[int:SPITRANSDONEINT]{SPI\_TRANS\_DONE\_INT} 的中断清除位。 (WT)
\label{fielddesc:SPIDMASEGTRANSDONEINTCLR}\item [SPI\_DMA\_SEG\_TRANS\_DONE\_INT\_CLR] \hyperref[int:SPIDMASEGTRANSDONEINT]{SPI\_DMA\_SEG\_TRANS\_DONE\_INT} 的中断清除位。 (WT)
\label{fielddesc:SPISEGMAGICERRINTCLR}\item [\textcolor{red}{SPI\_SEG\_MAGIC\_ERR\_INT\_CLR（仅对 SPI2 有效）}] \hyperref[int:SPISEGMAGICERRINT]{SPI\_SEG\_MAGIC\_ERR\_INT} 的中断清除位。 (WT)
\label{fielddesc:SPISLVCMDERRINTCLR}\item [SPI\_SLV\_CMD\_ERR\_INT\_CLR] \hyperref[int:SPISLVCMDERRINT]{SPI\_SLV\_CMD\_ERR\_INT} 的中断清除位。 (WT)

\item [见下页]
\end{reglist}\end{regdesc}
\end{register}

\addtocounter{Regfloat}{-1}
\begin{register}{H}{\textcolor{red}{SPI\_DMA\_INT\_CLR\_REG}}{0x{}0038}
\begin{regdesc}\begin{reglist}
\item [接上页]
\label{fielddesc:SPIMSTRXAFIFOWFULLERRINTCLR}\item [SPI\_MST\_RX\_AFIFO\_WFULL\_ERR\_INT\_CLR] \hyperref[int:SPIMSTRXAFIFOWFULLERRINT]{SPI\_MST\_RX\_AFIFO\_WFULL\_ERR\_INT} 的中断清除位。 (WT)
\label{fielddesc:SPIMSTTXAFIFOREMPTYERRINTCLR}\item [SPI\_MST\_TX\_AFIFO\_REMPTY\_ERR\_INT\_CLR] \hyperref[int:SPIMSTTXAFIFOREMPTYERRINT]{SPI\_MST\_TX\_AFIFO\_REMPTY\_ERR\_INT} 的中断清除位。 (WT)
\label{fielddesc:SPIAPP2INTCLR}\item [SPI\_APP2\_INT\_CLR] \hyperref[int:SPIAPP2INT]{SPI\_APP2\_INT} 的中断清除位。 (WT)
\label{fielddesc:SPIAPP1INTCLR}\item [SPI\_APP1\_INT\_CLR] \hyperref[int:SPIAPP1INT]{SPI\_APP1\_INT} 的中断清除位。 (WT)
\end{reglist}\end{regdesc}
\end{register}


\begin{register}{H}{\textcolor{red}{SPI\_DMA\_INT\_RAW\_REG}}{0x{}003C}\label{regdesc:SPIDMAINTRAWREG}
\regfield{(reserved)}{11}{21}{00000000000}%
\regfield{SPI\_APP1\_INT\_RAW}{1}{20}{{0}}%
\regfield{SPI\_APP2\_INT\_RAW}{1}{19}{{0}}%
\regfield{SPI\_MST\_TX\_AFIFO\_REMPTY\_ERR\_INT\_RAW}{1}{18}{{0}}%
\regfield{SPI\_MST\_RX\_AFIFO\_WFULL\_ERR\_INT\_RAW}{1}{17}{{0}}%
\regfield{SPI\_SLV\_CMD\_ERR\_INT\_RAW}{1}{16}{{0}}%
\regfield{(reserved)}{1}{15}{{0}}%
\regfield{(reserved)|\textcolor{red}{SPI\_SEG\_MAGIC\_ERR\_INT\_RAW}}{1}{14}{{0}}%
\regfield{SPI\_DMA\_SEG\_TRANS\_DONE\_INT\_RAW}{1}{13}{{0}}%
\regfield{SPI\_TRANS\_DONE\_INT\_RAW}{1}{12}{{0}}%
\regfield{SPI\_SLV\_WR\_BUF\_DONE\_INT\_RAW}{1}{11}{{0}}%
\regfield{SPI\_SLV\_RD\_BUF\_DONE\_INT\_RAW}{1}{10}{{0}}%
\regfield{SPI\_SLV\_WR\_DMA\_DONE\_INT\_RAW}{1}{9}{{0}}%
\regfield{SPI\_SLV\_RD\_DMA\_DONE\_INT\_RAW}{1}{8}{{0}}%
\regfield{SPI\_SLV\_CMDA\_INT\_RAW}{1}{7}{{0}}%
\regfield{SPI\_SLV\_CMD9\_INT\_RAW}{1}{6}{{0}}%
\regfield{SPI\_SLV\_CMD8\_INT\_RAW}{1}{5}{{0}}%
\regfield{SPI\_SLV\_CMD7\_INT\_RAW}{1}{4}{{0}}%
\regfield{SPI\_SLV\_EN\_QPI\_INT\_RAW}{1}{3}{{0}}%
\regfield{SPI\_SLV\_EX\_QPI\_INT\_RAW}{1}{2}{{0}}%
\regfield{SPI\_DMA\_OUTFIFO\_EMPTY\_ERR\_INT\_RAW}{1}{1}{{0}}%
\regfield{SPI\_DMA\_INFIFO\_FULL\_ERR\_INT\_RAW}{1}{0}{{0}}%
\reglabel{Reset}\regnewline%
\begin{regdesc}\begin{reglist}
\label{fielddesc:SPIDMAINFIFOFULLERRINTRAW}\item [SPI\_DMA\_INFIFO\_FULL\_ERR\_INT\_RAW] \hyperref[int:SPIINFIFOFULLERRINT]{SPI\_DMA\_INFIFO\_FULL\_ERR\_INT} 的原始中断位。 (R/WTC/SS)
\label{fielddesc:SPIDMAOUTFIFOEMPTYERRINTRAW}\item [SPI\_DMA\_OUTFIFO\_EMPTY\_ERR\_INT\_RAW] \hyperref[int:SPIOUTFIFOEMPTYERRINT]{SPI\_DMA\_OUTFIFO\_EMPTY\_ERR\_INT} 的原始中断位。 (R/WTC/SS)
\label{fielddesc:SPISLVEXQPIINTRAW}\item [SPI\_SLV\_EX\_QPI\_INT\_RAW] \hyperref[int:SPISLVEXQPIINT]{SPI\_SLV\_EX\_QPI\_INT} 的原始中断位。 (R/WTC/SS)
\label{fielddesc:SPISLVENQPIINTRAW}\item [SPI\_SLV\_EN\_QPI\_INT\_RAW] \hyperref[int:SPISLVENQPIINT]{SPI\_SLV\_EN\_QPI\_INT} 的原始中断位。 (R/WTC/SS)
\label{fielddesc:SPISLVCMD7INTRAW}\item [SPI\_SLV\_CMD7\_INT\_RAW] \hyperref[int:SPISLVCMD7INT]{SPI\_SLV\_CMD7\_INT} 的原始中断位。 (R/WTC/SS)
\label{fielddesc:SPISLVCMD8INTRAW}\item [SPI\_SLV\_CMD8\_INT\_RAW] \hyperref[int:SPISLVCMD8INT]{SPI\_SLV\_CMD8\_INT} 的原始中断位。 (R/WTC/SS)
\label{fielddesc:SPISLVCMD9INTRAW}\item [SPI\_SLV\_CMD9\_INT\_RAW] \hyperref[int:SPISLVCMD9INT]{SPI\_SLV\_CMD9\_INT} 的原始中断位。 (R/WTC/SS)
\label{fielddesc:SPISLVCMDAINTRAW}\item [SPI\_SLV\_CMDA\_INT\_RAW] \hyperref[int:SPISLVCMDAINT]{SPI\_SLV\_CMDA\_INT} 的原始中断位。 (R/WTC/SS)
\label{fielddesc:SPISLVRDDMADONEINTRAW}\item [SPI\_SLV\_RD\_DMA\_DONE\_INT\_RAW] \hyperref[int:SPISLVRDDMADONEINT]{SPI\_SLV\_RD\_DMA\_DONE\_INT} 的原始中断位。 (R/WTC/SS)
\label{fielddesc:SPISLVWRDMADONEINTRAW}\item [SPI\_SLV\_WR\_DMA\_DONE\_INT\_RAW] \hyperref[int:SPISLVWRDMADONEINT]{SPI\_SLV\_WR\_DMA\_DONE\_INT} 的原始中断位。 (R/WTC/SS)
\label{fielddesc:SPISLVRDBUFDONEINTRAW}\item [SPI\_SLV\_RD\_BUF\_DONE\_INT\_RAW] \hyperref[int:SPISLVRDBUFDONEINT]{SPI\_SLV\_RD\_BUF\_DONE\_INT} 的原始中断位。 (R/WTC/SS)
\label{fielddesc:SPISLVWRBUFDONEINTRAW}\item [SPI\_SLV\_WR\_BUF\_DONE\_INT\_RAW] \hyperref[int:SPISLVWRBUFDONEINT]{SPI\_SLV\_WR\_BUF\_DONE\_INT} 的原始中断位。 (R/WTC/SS)
\label{fielddesc:SPITRANSDONEINTRAW}\item [SPI\_TRANS\_DONE\_INT\_RAW] \hyperref[int:SPITRANSDONEINT]{SPI\_TRANS\_DONE\_INT} 的原始中断位。 (R/WTC/SS)

\item [见下页]
\end{reglist}\end{regdesc}
\end{register}

\addtocounter{Regfloat}{-1}
\begin{register}{H}{\textcolor{red}{SPI\_DMA\_INT\_RAW\_REG}}{0x{}003C}
\begin{regdesc}\begin{reglist}
\item [接上页]
\label{fielddesc:SPIDMASEGTRANSDONEINTRAW}\item [SPI\_DMA\_SEG\_TRANS\_DONE\_INT\_RAW] \hyperref[int:SPIDMASEGTRANSDONEINT]{SPI\_DMA\_SEG\_TRANS\_DONE\_INT} 的原始中断位。 (R/WTC/SS)
\label{fielddesc:SPISEGMAGICERRINTRAW}\item [\textcolor{red}{SPI\_SEG\_MAGIC\_ERR\_INT\_RAW（仅对 SPI2 有效）}] \hyperref[int:SPISEGMAGICERRINT]{SPI\_SEG\_MAGIC\_ERR\_INT} 的原始中断位。 (R/WTC/SS)
\label{fielddesc:SPISLVCMDERRINTRAW}\item [SPI\_SLV\_CMD\_ERR\_INT\_RAW] \hyperref[int:SPISLVCMDERRINT]{SPI\_SLV\_CMD\_ERR\_INT} 的原始中断位。 (R/WTC/SS)
\label{fielddesc:SPIMSTRXAFIFOWFULLERRINTRAW}\item [SPI\_MST\_RX\_AFIFO\_WFULL\_ERR\_INT\_RAW] \hyperref[int:SPIMSTRXAFIFOWFULLERRINT]{SPI\_MST\_RX\_AFIFO\_WFULL\_ERR\_INT} 的原始中断位。 (R/WTC/SS)
\label{fielddesc:SPIMSTTXAFIFOREMPTYERRINTRAW}\item [SPI\_MST\_TX\_AFIFO\_REMPTY\_ERR\_INT\_RAW] \hyperref[int:SPIMSTTXAFIFOREMPTYERRINT]{SPI\_MST\_TX\_AFIFO\_REMPTY\_ERR\_INT} 的原始中断位。 (R/WTC/SS)
\label{fielddesc:SPIAPP2INTRAW}\item [SPI\_APP2\_INT\_RAW] \hyperref[int:SPIAPP2INT]{SPI\_APP2\_INT} 的原始中断位。该值仅由软件控制。 (R/WTC/SS)
\label{fielddesc:SPIAPP1INTRAW}\item [SPI\_APP1\_INT\_RAW] \hyperref[int:SPIAPP1INT]{SPI\_APP1\_INT} 的中断使能位。该值仅由软件控制。 (R/WTC/SS)
\end{reglist}\end{regdesc}
\end{register}


\begin{register}{H}{\textcolor{red}{SPI\_DMA\_INT\_ST\_REG}}{0x{}0040}\label{regdesc:SPIDMAINTSTREG}
\regfield{(reserved)}{11}{21}{00000000000}%
\regfield{SPI\_APP1\_INT\_ST}{1}{20}{{0}}%
\regfield{SPI\_APP2\_INT\_ST}{1}{19}{{0}}%
\regfield{SPI\_MST\_TX\_AFIFO\_REMPTY\_ERR\_INT\_ST}{1}{18}{{0}}%
\regfield{SPI\_MST\_RX\_AFIFO\_WFULL\_ERR\_INT\_ST}{1}{17}{{0}}%
\regfield{SPI\_SLV\_CMD\_ERR\_INT\_ST}{1}{16}{{0}}%
\regfield{(reserved)}{1}{15}{{0}}%
\regfield{(reserved)|\textcolor{red}{SPI\_SEG\_MAGIC\_ERR\_INT\_ST}}{1}{14}{{0}}%
\regfield{SPI\_DMA\_SEG\_TRANS\_DONE\_INT\_ST}{1}{13}{{0}}%
\regfield{SPI\_TRANS\_DONE\_INT\_ST}{1}{12}{{0}}%
\regfield{SPI\_SLV\_WR\_BUF\_DONE\_INT\_ST}{1}{11}{{0}}%
\regfield{SPI\_SLV\_RD\_BUF\_DONE\_INT\_ST}{1}{10}{{0}}%
\regfield{SPI\_SLV\_WR\_DMA\_DONE\_INT\_ST}{1}{9}{{0}}%
\regfield{SPI\_SLV\_RD\_DMA\_DONE\_INT\_ST}{1}{8}{{0}}%
\regfield{SPI\_SLV\_CMDA\_INT\_ST}{1}{7}{{0}}%
\regfield{SPI\_SLV\_CMD9\_INT\_ST}{1}{6}{{0}}%
\regfield{SPI\_SLV\_CMD8\_INT\_ST}{1}{5}{{0}}%
\regfield{SPI\_SLV\_CMD7\_INT\_ST}{1}{4}{{0}}%
\regfield{SPI\_SLV\_EN\_QPI\_INT\_ST}{1}{3}{{0}}%
\regfield{SPI\_SLV\_EX\_QPI\_INT\_ST}{1}{2}{{0}}%
\regfield{SPI\_DMA\_OUTFIFO\_EMPTY\_ERR\_INT\_ST}{1}{1}{{0}}%
\regfield{SPI\_DMA\_INFIFO\_FULL\_ERR\_INT\_ST}{1}{0}{{0}}%
\reglabel{Reset}\regnewline%
\begin{regdesc}\begin{reglist}
\label{fielddesc:SPIDMAINFIFOFULLERRINTST}\item [SPI\_DMA\_INFIFO\_FULL\_ERR\_INT\_ST] \hyperref[int:SPIINFIFOFULLERRINT]{SPI\_DMA\_INFIFO\_FULL\_ERR\_INT} 的中断状态位。 (RO)
\label{fielddesc:SPIDMAOUTFIFOEMPTYERRINTST}\item [SPI\_DMA\_OUTFIFO\_EMPTY\_ERR\_INT\_ST] \hyperref[int:SPIOUTFIFOEMPTYERRINT]{SPI\_DMA\_OUTFIFO\_EMPTY\_ERR\_INT} 的中断状态位。 (RO)
\label{fielddesc:SPISLVEXQPIINTST}\item [SPI\_SLV\_EX\_QPI\_INT\_ST] \hyperref[int:SPISLVEXQPIINT]{SPI\_SLV\_EX\_QPI\_INT} 的中断状态位。 (RO)
\label{fielddesc:SPISLVENQPIINTST}\item [SPI\_SLV\_EN\_QPI\_INT\_ST] \hyperref[int:SPISLVENQPIINT]{SPI\_SLV\_EN\_QPI\_INT} 的中断状态位。 (RO)
\label{fielddesc:SPISLVCMD7INTST}\item [SPI\_SLV\_CMD7\_INT\_ST] \hyperref[int:SPISLVCMD7INT]{SPI\_SLV\_CMD7\_INT} 的中断状态位。 (RO)
\label{fielddesc:SPISLVCMD8INTST}\item [SPI\_SLV\_CMD8\_INT\_ST] \hyperref[int:SPISLVCMD8INT]{SPI\_SLV\_CMD8\_INT} 的中断状态位。 (RO)
\label{fielddesc:SPISLVCMD9INTST}\item [SPI\_SLV\_CMD9\_INT\_ST] \hyperref[int:SPISLVCMD9INT]{SPI\_SLV\_CMD9\_INT} 的中断状态位。 (RO)
\label{fielddesc:SPISLVCMDAINTST}\item [SPI\_SLV\_CMDA\_INT\_ST] \hyperref[int:SPISLVCMDAINT]{SPI\_SLV\_CMDA\_INT} 的中断状态位。 (RO)
\label{fielddesc:SPISLVRDDMADONEINTST}\item [SPI\_SLV\_RD\_DMA\_DONE\_INT\_ST] \hyperref[int:SPISLVRDDMADONEINT]{SPI\_SLV\_RD\_DMA\_DONE\_INT} 的中断状态位。 (RO)
\label{fielddesc:SPISLVWRDMADONEINTST}\item [SPI\_SLV\_WR\_DMA\_DONE\_INT\_ST] \hyperref[int:SPISLVWRDMADONEINT]{SPI\_SLV\_WR\_DMA\_DONE\_INT} 的中断状态位。 (RO)
\label{fielddesc:SPISLVRDBUFDONEINTST}\item [SPI\_SLV\_RD\_BUF\_DONE\_INT\_ST] \hyperref[int:SPISLVRDBUFDONEINT]{SPI\_SLV\_RD\_BUF\_DONE\_INT} 的中断状态位。 (RO)
\label{fielddesc:SPISLVWRBUFDONEINTST}\item [SPI\_SLV\_WR\_BUF\_DONE\_INT\_ST] \hyperref[int:SPISLVWRBUFDONEINT]{SPI\_SLV\_WR\_BUF\_DONE\_INT} 的中断状态位。 (RO)
\label{fielddesc:SPITRANSDONEINTST}\item [SPI\_TRANS\_DONE\_INT\_ST] \hyperref[int:SPITRANSDONEINT]{SPI\_TRANS\_DONE\_INT} 的中断状态位。 (RO)
\label{fielddesc:SPIDMASEGTRANSDONEINTST}\item [SPI\_DMA\_SEG\_TRANS\_DONE\_INT\_ST] \hyperref[int:SPIDMASEGTRANSDONEINT]{SPI\_DMA\_SEG\_TRANS\_DONE\_INT} 的中断状态位。 (RO)
\label{fielddesc:SPISEGMAGICERRINTST}\item [\textcolor{red}{SPI\_SEG\_MAGIC\_ERR\_INT\_ST（仅对 SPI2 有效）}] \hyperref[int:SPISEGMAGICERRINT]{SPI\_SEG\_MAGIC\_ERR\_INT} 的中断状态位。 (RO)
\label{fielddesc:SPISLVCMDERRINTST}\item [SPI\_SLV\_CMD\_ERR\_INT\_ST] \hyperref[int:SPISLVCMDERRINT]{SPI\_SLV\_CMD\_ERR\_INT} 的中断状态位。 (RO)
\label{fielddesc:SPIMSTRXAFIFOWFULLERRINTST}\item [SPI\_MST\_RX\_AFIFO\_WFULL\_ERR\_INT\_ST] \hyperref[int:SPIMSTRXAFIFOWFULLERRINT]{SPI\_MST\_RX\_AFIFO\_WFULL\_ERR\_INT} 的中断状态位。 (RO)


\item [见下页]
\end{reglist}\end{regdesc}
\end{register}

\addtocounter{Regfloat}{-1}
\begin{register}{H}{\textcolor{red}{SPI\_DMA\_INT\_ST\_REG}}{0x{}0040}
\begin{regdesc}\begin{reglist}
\item [接上页]
\label{fielddesc:SPIMSTTXAFIFOREMPTYERRINTST}\item [SPI\_MST\_TX\_AFIFO\_REMPTY\_ERR\_INT\_ST] \hyperref[int:SPIMSTTXAFIFOREMPTYERRINT]{SPI\_MST\_TX\_AFIFO\_REMPTY\_ERR\_INT} 的中断状态位。 (RO)
\label{fielddesc:SPIAPP2INTST}\item [SPI\_APP2\_INT\_ST] \hyperref[int:SPIAPP2INT]{SPI\_APP2\_INT} 的中断状态位。 (RO)
\label{fielddesc:SPIAPP1INTST}\item [SPI\_APP1\_INT\_ST] \hyperref[int:SPIAPP1INT]{SPI\_APP1\_INT} 的中断状态位。 (RO)
\end{reglist}\end{regdesc}
\end{register}


\begin{register}{H}{\textcolor{red}{SPI\_DMA\_INT\_SET\_REG}}{0x{}0044}\label{regdesc:SPIDMAINTSETREG}
\regfield{(reserved)}{11}{21}{00000000000}%
\regfield{SPI\_APP1\_INT\_SET}{1}{20}{{0}}%
\regfield{SPI\_APP2\_INT\_SET}{1}{19}{{0}}%
\regfield{SPI\_MST\_TX\_AFIFO\_REMPTY\_ERR\_INT\_SET}{1}{18}{{0}}%
\regfield{SPI\_MST\_RX\_AFIFO\_WFULL\_ERR\_INT\_SET}{1}{17}{{0}}%
\regfield{SPI\_SLV\_CMD\_ERR\_INT\_SET}{1}{16}{{0}}%
\regfield{(reserved)}{1}{15}{{0}}%
\regfield{(reserved)|\textcolor{red}{SPI\_SEG\_MAGIC\_ERR\_INT\_SET}}{1}{14}{{0}}%
\regfield{SPI\_DMA\_SEG\_TRANS\_DONE\_INT\_SET}{1}{13}{{0}}%
\regfield{SPI\_TRANS\_DONE\_INT\_SET}{1}{12}{{0}}%
\regfield{SPI\_SLV\_WR\_BUF\_DONE\_INT\_SET}{1}{11}{{0}}%
\regfield{SPI\_SLV\_RD\_BUF\_DONE\_INT\_SET}{1}{10}{{0}}%
\regfield{SPI\_SLV\_WR\_DMA\_DONE\_INT\_SET}{1}{9}{{0}}%
\regfield{SPI\_SLV\_RD\_DMA\_DONE\_INT\_SET}{1}{8}{{0}}%
\regfield{SPI\_SLV\_CMDA\_INT\_SET}{1}{7}{{0}}%
\regfield{SPI\_SLV\_CMD9\_INT\_SET}{1}{6}{{0}}%
\regfield{SPI\_SLV\_CMD8\_INT\_SET}{1}{5}{{0}}%
\regfield{SPI\_SLV\_CMD7\_INT\_SET}{1}{4}{{0}}%
\regfield{SPI\_SLV\_EN\_QPI\_INT\_SET}{1}{3}{{0}}%
\regfield{SPI\_SLV\_EX\_QPI\_INT\_SET}{1}{2}{{0}}%
\regfield{SPI\_DMA\_OUTFIFO\_EMPTY\_ERR\_INT\_SET}{1}{1}{{0}}%
\regfield{SPI\_DMA\_INFIFO\_FULL\_ERR\_INT\_SET}{1}{0}{{0}}%
\reglabel{Reset}\regnewline%
\begin{regdesc}\begin{reglist}
\label{fielddesc:SPIDMAINFIFOFULLERRINTSET}\item [SPI\_DMA\_INFIFO\_FULL\_ERR\_INT\_SET] 软件置位 \hyperref[int:SPIINFIFOFULLERRINT]{SPI\_DMA\_INFIFO\_FULL\_ERR\_INT} 中断。 (WT)
\label{fielddesc:SPIDMAOUTFIFOEMPTYERRINTSET}\item [SPI\_DMA\_OUTFIFO\_EMPTY\_ERR\_INT\_SET] 软件置位\hyperref[int:SPIOUTFIFOEMPTYERRINT]{SPI\_DMA\_OUTFIFO\_EMPTY\_ERR\_INT} 中断。 (WT)
\label{fielddesc:SPISLVEXQPIINTSET}\item [SPI\_SLV\_EX\_QPI\_INT\_SET] 软件置位 \hyperref[int:SPISLVEXQPIINT]{SPI\_SLV\_EX\_QPI\_INT} 中断。 (WT)
\label{fielddesc:SPISLVENQPIINTSET}\item [SPI\_SLV\_EN\_QPI\_INT\_SET] 软件置位 \hyperref[int:SPISLVENQPIINT]{SPI\_SLV\_EN\_QPI\_INT} 中断。 (WT)
\label{fielddesc:SPISLVCMD7INTSET}\item [SPI\_SLV\_CMD7\_INT\_SET] 软件置位 \hyperref[int:SPISLVCMD7INT]{SPI\_SLV\_CMD7\_INT} 中断。 (WT)
\label{fielddesc:SPISLVCMD8INTSET}\item [SPI\_SLV\_CMD8\_INT\_SET] 软件置位 \hyperref[int:SPISLVCMD8INT]{SPI\_SLV\_CMD8\_INT} 中断。 (WT)
\label{fielddesc:SPISLVCMD9INTSET}\item [SPI\_SLV\_CMD9\_INT\_SET] 软件置位 \hyperref[int:SPISLVCMD9INT]{SPI\_SLV\_CMD9\_INT} 中断。 (WT)
\label{fielddesc:SPISLVCMDAINTSET}\item [SPI\_SLV\_CMDA\_INT\_SET] 软件置位 \hyperref[int:SPISLVCMDAINT]{SPI\_SLV\_CMDA\_INT} 中断。 (WT)

\item [见下页]
\end{reglist}\end{regdesc}
\end{register}

\addtocounter{Regfloat}{-1}
\begin{register}{H}{\textcolor{red}{SPI\_DMA\_INT\_SET\_REG}}{0x{}0044}
\begin{regdesc}\begin{reglist}
\item [接上页]
\label{fielddesc:SPISLVRDDMADONEINTSET}\item [SPI\_SLV\_RD\_DMA\_DONE\_INT\_SET] 软件置位 \hyperref[int:SPISLVRDDMADONEINT]{SPI\_SLV\_RD\_DMA\_DONE\_INT} 中断。 (WT)
\label{fielddesc:SPISLVWRDMADONEINTSET}\item [SPI\_SLV\_WR\_DMA\_DONE\_INT\_SET] 软件置位 \hyperref[int:SPISLVWRDMADONEINT]{SPI\_SLV\_WR\_DMA\_DONE\_INT} 中断。 (WT)
\label{fielddesc:SPISLVRDBUFDONEINTSET}\item [SPI\_SLV\_RD\_BUF\_DONE\_INT\_SET] 软件置位 \hyperref[int:SPISLVRDBUFDONEINT]{SPI\_SLV\_RD\_BUF\_DONE\_INT} 中断。 (WT)
\label{fielddesc:SPISLVWRBUFDONEINTSET}\item [SPI\_SLV\_WR\_BUF\_DONE\_INT\_SET] 软件置位 \hyperref[int:SPISLVWRBUFDONEINT]{SPI\_SLV\_WR\_BUF\_DONE\_INT} 中断。 (WT)
\label{fielddesc:SPITRANSDONEINTSET}\item [SPI\_TRANS\_DONE\_INT\_SET] 软件置位 \hyperref[int:SPITRANSDONEINT]{SPI\_TRANS\_DONE\_INT} 中断。 (WT)
\label{fielddesc:SPIDMASEGTRANSDONEINTSET}\item [SPI\_DMA\_SEG\_TRANS\_DONE\_INT\_SET] 软件置位 \hyperref[int:SPIDMASEGTRANSDONEINT]{SPI\_DMA\_SEG\_TRANS\_DONE\_INT} 中断。 (WT)
\label{fielddesc:SPISEGMAGICERRINTSET}\item [\textcolor{red}{SPI\_SEG\_MAGIC\_ERR\_INT\_SET（仅对 SPI2 有效）}] 软件置位 \hyperref[int:SPISEGMAGICERRINT]{SPI\_SEG\_MAGIC\_ERR\_INT} 中断。 (WT)
\label{fielddesc:SPISLVCMDERRINTSET}\item [SPI\_SLV\_CMD\_ERR\_INT\_SET] 软件置位 \hyperref[int:SPISLVCMDERRINT]{SPI\_SLV\_CMD\_ERR\_INT} 中断。 (WT)
\label{fielddesc:SPIMSTRXAFIFOWFULLERRINTSET}\item [SPI\_MST\_RX\_AFIFO\_WFULL\_ERR\_INT\_SET] 软件置位 \hyperref[int:SPIMSTRXAFIFOWFULLERRINT]{SPI\_MST\_RX\_AFIFO\_WFULL\_ERR\_INT} 中断。 (WT)
\label{fielddesc:SPIMSTTXAFIFOREMPTYERRINTSET}\item [SPI\_MST\_TX\_AFIFO\_REMPTY\_ERR\_INT\_SET] 软件置位 \hyperref[int:SPIMSTTXAFIFOREMPTYERRINT]{SPI\_MST\_TX\_AFIFO\_REMPTY\_ERR\_INT} 中断。 (WT)
\label{fielddesc:SPIAPP2INTSET}\item [SPI\_APP2\_INT\_SET] 软件置位 \hyperref[int:SPIAPP2INT]{SPI\_APP2\_INT} 中断。 (WT)
\label{fielddesc:SPIAPP1INTSET}\item [SPI\_APP1\_INT\_SET] 软件置位 \hyperref[int:SPIAPP1INT]{SPI\_APP1\_INT} 中断。 (WT)
\end{reglist}\end{regdesc}
\end{register}


\begin{register}{H}{SPI\_W0\_REG}{0x{}0098}\label{regdesc:SPIW0REG}
\regfield{SPI\_BUF0}{32}{0}{{0}}%
\reglabel{Reset}\regnewline%
\begin{regdesc}\begin{reglist}
\label{fielddesc:SPIBUF0}\item [SPI\_BUF0] 数据 buffer 0，32 位。 (R/W/SS)
\end{reglist}\end{regdesc}
\end{register}


\begin{register}{H}{SPI\_W1\_REG}{0x{}009C}\label{regdesc:SPIW1REG}
\regfield{SPI\_BUF1}{32}{0}{{0}}%
\reglabel{Reset}\regnewline%
\begin{regdesc}\begin{reglist}
\label{fielddesc:SPIBUF1}\item [SPI\_BUF1] 数据 buffer 1，32 位。 (R/W/SS)
\end{reglist}\end{regdesc}
\end{register}


\begin{register}{H}{SPI\_W2\_REG}{0x{}00A0}\label{regdesc:SPIW2REG}
\regfield{SPI\_BUF2}{32}{0}{{0}}%
\reglabel{Reset}\regnewline%
\begin{regdesc}\begin{reglist}
\label{fielddesc:SPIBUF2}\item [SPI\_BUF2] 数据 buffer 2，32 位。 (R/W/SS)
\end{reglist}\end{regdesc}
\end{register}


\begin{register}{H}{SPI\_W3\_REG}{0x{}00A4}\label{regdesc:SPIW3REG}
\regfield{SPI\_BUF3}{32}{0}{{0}}%
\reglabel{Reset}\regnewline%
\begin{regdesc}\begin{reglist}
\label{fielddesc:SPIBUF3}\item [SPI\_BUF3] 数据 buffer 3，32 位。 (R/W/SS)
\end{reglist}\end{regdesc}
\end{register}


\begin{register}{H}{SPI\_W4\_REG}{0x{}00A8}\label{regdesc:SPIW4REG}
\regfield{SPI\_BUF4}{32}{0}{{0}}%
\reglabel{Reset}\regnewline%
\begin{regdesc}\begin{reglist}
\label{fielddesc:SPIBUF4}\item [SPI\_BUF4] 数据 buffer 4，32 位。 (R/W/SS)
\end{reglist}\end{regdesc}
\end{register}


\begin{register}{H}{SPI\_W5\_REG}{0x{}00AC}\label{regdesc:SPIW5REG}
\regfield{SPI\_BUF5}{32}{0}{{0}}%
\reglabel{Reset}\regnewline%
\begin{regdesc}\begin{reglist}
\label{fielddesc:SPIBUF5}\item [SPI\_BUF5] 数据 buffer 5，32 位。 (R/W/SS)
\end{reglist}\end{regdesc}
\end{register}


\begin{register}{H}{SPI\_W6\_REG}{0x{}00B0}\label{regdesc:SPIW6REG}
\regfield{SPI\_BUF6}{32}{0}{{0}}%
\reglabel{Reset}\regnewline%
\begin{regdesc}\begin{reglist}
\label{fielddesc:SPIBUF6}\item [SPI\_BUF6] 数据 buffer 6，32 位。 (R/W/SS)
\end{reglist}\end{regdesc}
\end{register}


\begin{register}{H}{SPI\_W7\_REG}{0x{}00B4}\label{regdesc:SPIW7REG}
\regfield{SPI\_BUF7}{32}{0}{{0}}%
\reglabel{Reset}\regnewline%
\begin{regdesc}\begin{reglist}
\label{fielddesc:SPIBUF7}\item [SPI\_BUF7] 数据 buffer 7，32 位。 (R/W/SS)
\end{reglist}\end{regdesc}
\end{register}


\begin{register}{H}{SPI\_W8\_REG}{0x{}00B8}\label{regdesc:SPIW8REG}
\regfield{SPI\_BUF8}{32}{0}{{0}}%
\reglabel{Reset}\regnewline%
\begin{regdesc}\begin{reglist}
\label{fielddesc:SPIBUF8}\item [SPI\_BUF8] 数据 buffer 8，32 位。 (R/W/SS)
\end{reglist}\end{regdesc}
\end{register}


\begin{register}{H}{SPI\_W9\_REG}{0x{}00BC}\label{regdesc:SPIW9REG}
\regfield{SPI\_BUF9}{32}{0}{{0}}%
\reglabel{Reset}\regnewline%
\begin{regdesc}\begin{reglist}
\label{fielddesc:SPIBUF9}\item [SPI\_BUF9] 数据 buffer 9，32 位。 (R/W/SS)
\end{reglist}\end{regdesc}
\end{register}


\begin{register}{H}{SPI\_W10\_REG}{0x{}00C0}\label{regdesc:SPIW10REG}
\regfield{SPI\_BUF10}{32}{0}{{0}}%
\reglabel{Reset}\regnewline%
\begin{regdesc}\begin{reglist}
\label{fielddesc:SPIBUF10}\item [SPI\_BUF10] 数据 buffer 10，32 位。 (R/W/SS)
\end{reglist}\end{regdesc}
\end{register}


\begin{register}{H}{SPI\_W11\_REG}{0x{}00C4}\label{regdesc:SPIW11REG}
\regfield{SPI\_BUF11}{32}{0}{{0}}%
\reglabel{Reset}\regnewline%
\begin{regdesc}\begin{reglist}
\label{fielddesc:SPIBUF11}\item [SPI\_BUF11] 数据 buffer 11，32 位。 (R/W/SS)
\end{reglist}\end{regdesc}
\end{register}


\begin{register}{H}{SPI\_W12\_REG}{0x{}00C8}\label{regdesc:SPIW12REG}
\regfield{SPI\_BUF12}{32}{0}{{0}}%
\reglabel{Reset}\regnewline%
\begin{regdesc}\begin{reglist}
\label{fielddesc:SPIBUF12}\item [SPI\_BUF12] 数据 buffer 12，32 位。 (R/W/SS)
\end{reglist}\end{regdesc}
\end{register}


\begin{register}{H}{SPI\_W13\_REG}{0x{}00CC}\label{regdesc:SPIW13REG}
\regfield{SPI\_BUF13}{32}{0}{{0}}%
\reglabel{Reset}\regnewline%
\begin{regdesc}\begin{reglist}
\label{fielddesc:SPIBUF13}\item [SPI\_BUF13] 数据 buffer 13，32 位。 (R/W/SS)
\end{reglist}\end{regdesc}
\end{register}


\begin{register}{H}{SPI\_W14\_REG}{0x{}00D0}\label{regdesc:SPIW14REG}
\regfield{SPI\_BUF14}{32}{0}{{0}}%
\reglabel{Reset}\regnewline%
\begin{regdesc}\begin{reglist}
\label{fielddesc:SPIBUF14}\item [SPI\_BUF14] 数据 buffer 14，32 位。 (R/W/SS)
\end{reglist}\end{regdesc}
\end{register}


\begin{register}{H}{SPI\_W15\_REG}{0x{}00D4}\label{regdesc:SPIW15REG}
\regfield{SPI\_BUF15}{32}{0}{{0}}%
\reglabel{Reset}\regnewline%
\begin{regdesc}\begin{reglist}
\label{fielddesc:SPIBUF15}\item [SPI\_BUF15] 数据 buffer 15，32 位。 (R/W/SS)
\end{reglist}\end{regdesc}
\end{register}


\begin{register}{H}{SPI\_DATE\_REG}{0x{}00F0}\label{regdesc:SPIDATEREG}
\regfield{(reserved)}{4}{28}{0000}%
\regfield{SPI\_DATE}{28}{0}{{0x{}2101190}}%
\reglabel{Reset}\regnewline%
\begin{regdesc}\begin{reglist}
\label{fielddesc:SPIDATE}\item [SPI\_DATE] 版本寄存器。 (R/W)
\end{reglist}\end{regdesc}
\end{register}


